\documentclass[11pt,reqno]{amsart}

\usepackage[a4paper,margin=28mm]{geometry}
\usepackage{amsmath,amssymb,mathtools,mathrsfs}
\usepackage{enumitem,booktabs,longtable,array}
\usepackage{microtype}
\usepackage{xurl}
\usepackage[colorlinks=true,linkcolor=blue,citecolor=blue,urlcolor=blue]{hyperref}

\numberwithin{equation}{section}
\setlist[enumerate]{label=\textup{(\roman*)},leftmargin=2.2em}

\newtheorem{theorem}{Theorem}[section]
\newtheorem{proposition}[theorem]{Proposition}
\newtheorem{lemma}[theorem]{Lemma}
\newtheorem{corollary}[theorem]{Corollary}
\theoremstyle{definition}
\newtheorem{example}[theorem]{Example}
\newtheorem{remark}[theorem]{Remark}
\theoremstyle{definition}
\newtheorem{definition}[theorem]{Definition}

\newcommand{\N}{\mathbb N}
\newcommand{\Nzero}{\mathbb N_0}
\newcommand{\C}{\mathcal C}
\newcommand{\D}{\mathcal D}
\newcommand{\Acal}{\mathcal A}
\newcommand{\Lcal}{\mathcal L}
\newcommand{\Patt}{\mathsf P}
\newcommand{\hseq}{h_{\mathrm{seq}}}
\newcommand{\hpat}{h^{*}}
\newcommand{\hinv}{h_{\mathrm{inv}}}
\newcommand{\htop}{h_{\mathrm{top}}}
\newcommand{\Bor}{\mathrm{Bor}}
\newcommand{\clop}{\mathrm{clop}}
\newcommand{\lang}{\mathrm{lang}}
\newcommand{\supp}{\operatorname{supp}}
\newcommand{\doi}[1]{\href{https://doi.org/#1}{doi:\nolinkurl{#1}}}


\newcommand{\K}{\mathcal K}
\newcommand{\Ssafe}{\mathcal S}
\newcommand{\hstar}{h^*}
\newcommand{\one}{\mathbf 1}
\newcommand{\opnorm}[1]{\left\|#1\right\|_{\infty\to\infty}}
\newcommand{\esssup}{\operatorname*{ess\,sup}}
\newcommand{\R}{\mathbb R}
\newcommand{\Leb}{\operatorname{Leb}}
\newcommand{\cl}{\operatorname{cl}}
\newcommand{\intr}{\operatorname{int}}
\newcommand{\sgn}{\operatorname{sgn}}
\newcommand{\shift}{\sigma}
\newcommand{\ind}{\operatorname{ind}}
\newcommand{\const}{\mathrm{const}}
\newcommand{\ellA}{\ell_2^{A}}
\newcommand{\eps}{\varepsilon}
\allowdisplaybreaks[1]
\hypersetup{pdftitle={Sparse Observations and Pattern Complexity of Safe Feedbacks},pdfauthor={Ziqing Ding}}
\title[Sparse observations of safe feedbacks]{Sparse Observations and Pattern Complexity of Safe Feedbacks}
\author{Ziqing Ding}
\address{School of Mathematical Sciences, Nanjing Normal University,
Nanjing 210023, China}
\date{}
\subjclass[2020]{93C55, 37B10, 37B40, 94A17}
\keywords{safe feedback, sequence entropy, maximal pattern entropy,
Borel partition, inverse capacity, symbolic dynamics, matrix products}
\begin{document}
\begin{abstract}
A safe state feedback generates a sequence of control instructions from
each initial state. We study the growth of the number of patterns seen
at selected times, with a fixed control period $\tau$. The main question
is whether one observation sequence detects complexity uniformly over
all safe Borel feedbacks. We answer it when all legal blocks have
a common endpoint map conjugate to a power of a mixing one-sided shift
of finite type. The minimum rate is $\log r/\tau$, where $r$ is the
minimum number of block safety domains covering the state space, and
every observation sequence with diverging gaps attains it. When controls
have different successors, we bound the initial-state sets producing
specified words by products of inverse-transfer operators. A contraction
condition, combined with the finite covering lemma of Kerr and Li,
determines the minimum maximal-pattern entropy; the observation times
in this argument may depend on the feedback. A separate deletion
argument gives common attaining sequences for countably many name
languages. Symbolic and affine examples distinguish these conclusions
from finite-horizon safe covering and from invariance entropy.
\end{abstract}
\maketitle

\section{Introduction}\label{sec:intro}

Consider a discrete-time control system and a compact set $Q$ in which
the state must remain. A feedback chooses a control from the current
state; the same rule is used at every subsequent step. Starting from
different points of $Q$ produces different sequences of controls. We
ask how few patterns these sequences can have when they are observed
at selected times. The feedback is allowed to be Borel, so the question
includes rules with irregular decision regions.

A simple symbolic system illustrates the choice involved. Let
$Q=\{0,1,2\}^{\Nzero}$. At a state $x$, either control different
from $x_0$ is safe, and both safe controls send $x$ to the shift
$\sigma x$. A feedback is therefore a Borel function $s$ satisfying
$s(x)\ne x_0$. Its control sequence is
\[
 s(x),\ s(\sigma x),\ s(\sigma^2x),\ldots.
\]
The state orbit is already determined, but the observed labels depend
on how $s$ resolves the overlaps of the three safety sets. Choosing
$s(x)=1$ when $x_0=0$ and $s(x)=0$ otherwise uses two labels and
realizes every binary pattern at any finite set of times. Can another
Borel choice reduce this complexity? If the observation times are
sufficiently far apart, our closed-cover theorem says that it cannot.
We return to this system in Sections~\ref{sec:selectors}
and~\ref{sec:closed-cover}.

More generally, fix a control period $\tau$: the chosen instruction is
a word of $\tau$ controls, all of whose intermediate states must stay
in $Q$. Write $\Patt_s(F)$ for the number of control-block patterns
produced by a feedback $s$ on a finite set $F$ of block indices. For
an increasing sequence $A=(a_i)$, its sequence entropy is
\[
 h_A(s)=\limsup_{n\to\infty}\frac{1}{n\tau}
       \log\Patt_s(\{a_1,\ldots,a_n\}).
\]
The normalization counts observations, each charged the block length
$\tau$; it does not count the intervening gaps. Maximizing over
$n$-element windows before taking the limit gives the maximal-pattern
entropy $h^*(s)$. Section~\ref{sec:control} gives the definitions
and the elementary limit argument.

There are two orders in which to choose a feedback and an observation
sequence:
\[
 L_\tau=\sup_A\inf_s h_A(s),\qquad
 M_\tau=\inf_s h^*(s).
\]
For each fixed $s$, one has $h^*(s)=\max_A h_A(s)$
(Corollary~\ref{cor:fixed}). Thus $M_\tau$ allows the observation
sequence to be chosen after the feedback. For $L_\tau$, the sequence
must be fixed first and give a lower bound for every safe feedback.
The inequality $L_\tau\le M_\tau$ is immediate; the opposite
inequality requires a way to choose observation times in common.

Invariance entropy addresses a related question at consecutive times.
It counts the smallest collection of control words that keeps every
initial state safe for a finite horizon, and measures growth per
elapsed time \cite{Kawan2013}. Such a collection can be chosen anew
for each horizon. A feedback, by contrast, must make compatible choices
whenever a trajectory returns to the same state. Sparse observations
also retain information which a consecutive-time rate can miss. These
differences matter even after optimization: the Thue--Morse realization
in Section~\ref{sec:thue-morse} has invariance entropy zero and
$L_1=M_1=\log2$. If the control period is allowed to vary, the
infimum of the pattern rates returns to invariance entropy
(Proposition~\ref{prop:all-period-collapse}).

We approach the fixed-period problem in two ways. First suppose that
every legal block has the same endpoint map $T$, and that $T$ is
conjugate to a power of the shift on a mixing one-sided shift of finite
type. Let $r$ be the minimum number of block safety domains covering
$Q$. Corollary~\ref{cc-thm:control} gives
\[
 \inf_s h_A(s)=L_\tau=M_\tau=\frac{\log r}{\tau}
 \quad\text{whenever }a_{i+1}-a_i\longrightarrow\infty.
\]
The upper bound uses an ordered minimum subcover. For the lower bound,
any refining Borel partition has at least $r$ nonmeager atoms. Each of
these atoms occupies most of some cylinder in the sense of Baire
category. Mixing joins sufficiently separated cylinders, and openness
of the shift allows the resulting orbit to avoid the meager exceptions.
This proves the bound on one prescribed sequence, simultaneously for
all Borel feedbacks.

When the successor depends on the chosen control, this concatenation
argument is no longer available. The second approach instead bounds
the size of the sets of initial states producing specified control
words. For a fixed list of $k-1$ allowed labels at each time, a product
of safe inverse-transfer operators bounds the mass carried by all words
in those lists. If these bounds decay exponentially, the actual names
cannot be covered by few such list boxes. The finite covering lemma
of Kerr and Li then supplies $k$ independent label choices on
arbitrarily large sets of coordinates. Theorem~\ref{thm:capacity-main}
deduces $h^*(s)\ge\log k/\tau$ for every feedback, and
$M_\tau=\log r/\tau$ when $k=r$. The independent positions may
depend on $s$, which is exactly the obstacle to computing $L_\tau$
by the same proof.

The combinatorial input is due to Kerr and Li
\cite[Lemma~3.3]{KL2007}; it also follows from Huang and Ye's
sharper estimates \cite[Lemma~4.1 and Remark~4.2]{HuangYe2009}.
The application here is the estimate for feedback fibers, with the
safety restrictions retained at every intermediate time. The reference
probability need not be invariant. Two examples distinguish the product
condition from simpler tests: separated density peaks defeat the scalar
test for a fixed measure, and a three-coordinate matrix product contracts
even though the single-step density test fails for every full-support
reference probability. Appendix~\ref{sec:cylinder} explains the exact
matrix reduction and gives a general sufficient condition using a
common positive left weight.

The closed-cover argument is related to Romagnoli's comparison of
an excluded-letter cover with its clopen refinements
\cite{Romagnoli2007}. Here the refinements may be Borel, and a
sparse sequence is prescribed before the refinement is chosen.
The measurable common-sequence results of Huang, Maass, and Ye
\cite[Lemma~3.3 and Theorem~3.5]{HuangMaassYe2004} concern
Shannon entropy relative to an invariant probability. Our counts
include every realized pattern, even one supported on exceptional
states. The broader sequence-entropy and independence framework is
developed in \cite{Goodman1974,Hulse1982,KamaeZamboniSystems2002,
KL2007,KL2009}.

The common-sequence question returns in Section~\ref{sec:comb}:
countably many name languages can share an attaining sequence even
without shift invariance. Section~\ref{sec:comparison} then explains
why fixing the control period is essential, using a forced symbolic
feedback and the Thue--Morse system. The affine examples in
Section~\ref{sec:affine} compare two ways to obtain a lower bound:
expansion makes individual fibers small, whereas protected inverse
branches realize all sampled words. Their difference remains visible
in a four-branch system for which $M_1$ is known but $L_1$ is not.
The appendices supply the finite combinatorial estimate, the matrix
criterion, the longer counts for this affine system, and a counterexample
to pointwise common attainment. Each can be read with the result in the
main text that calls for it.

\section{Safe feedbacks and their observed names}\label{sec:control}

\subsection{The control rule}\label{sec:selectors}
Let $X$ be a metrizable space and let $f_u:X\to X$ be continuous
for each control $u\in U$. Initially $U$ is finite; when we later
consider invariant partitions we also allow an arbitrary control set.
For a word $w=(w_0,\ldots,w_{N-1})$, define
\[
 \varphi(0,x,w)=x,\qquad
 \varphi(j+1,x,w)=f_{w_j}(\varphi(j,x,w)).
\]
Throughout, all control sequences in $U^{\Nzero}$ are admissible.
Let $Q\subset X$ be nonempty, compact and \emph{controlled invariant}:
from every $x\in Q$ some infinite control sequence stays in $Q$.
We use $\N=\{1,2,\ldots\}$, $\Nzero=\{0,1,\ldots\}$ and
natural logarithms.

Fix $\tau\in\N$ and write $W=U^\tau$. The safety domain of a
block $w$ and its endpoint map are
\begin{equation}\label{cc-eq:safe}
 K_w=\{x\in Q:\varphi(j,x,w)\in Q\ (0\le j\le\tau)\},
 \qquad g_w(x)=\varphi(\tau,x,w).
\end{equation}
Each $K_w$ is closed. These sets cover $Q$, since a safe infinite
control provides a safe first block. Their overlaps describe the
choices available to a feedback.

A \emph{safe feedback} is a Borel function $s:Q\to W$ such that
$x\in K_{s(x)}$ for every $x\in Q$. Its block-boundary map is
$T_s(x)=g_{s(x)}(x)$. This map is Borel and takes $Q$ into itself.
Applying the same rule at successive states keeps the entire trajectory
in $Q$, including the times between block boundaries. We denote the
class of all these rules by $\Ssafe_\tau$. The rule uses the current
state only: there is no additional clock or memory variable.
Continuity, or dependence on finitely many coordinates, is not required.

Sometimes we specify a finite block collection $W\subset U^\tau$
whose safety domains cover $Q$. In that case $\Ssafe_\tau$ and its
optimization values are relative to that collection. Unless such a
restriction is stated, $W=U^\tau$.

\subsection{Counting at selected times}
Starting at $x$, the rule produces the infinite name
$(s(T_s^jx))_{j\ge0}$. For a finite set $F\subset\Nzero$ put
\[
 \Patt_s(F)=\bigl|\{(s(T_s^jx))_{j\in F}:x\in Q\}\bigr|,
 \qquad \Patt_s(\varnothing)=1.
\]
Thus a pattern is counted if it occurs on at least one complete safe
trajectory. Write $\mathscr S$ for the strictly increasing infinite
sequences in $\Nzero$. For $A=(a_i)\in\mathscr S$, set
\[
 h_A(s)=\limsup_{n\to\infty}\frac{\log\Patt_s(\{a_1,\ldots,a_n\})}{n\tau},
 \qquad p_s^*(n)=\max_{|F|=n}\Patt_s(F),
\]
and
\[
 h^*(s)=\lim_{n\to\infty}\frac{\log p_s^*(n)}{n\tau}.
\]
The maximum exists because the possible counts are integers between
$1$ and $|W|^n$. Splitting a window into parts of sizes $m,n$ gives
$p_s^*(m+n)\le p_s^*(m)p_s^*(n)$, so Fekete's lemma gives the
last limit. In particular $h_A(s)\le h^*(s)$.

The index $a_i$ refers to the $a_i$th block, observed at physical
time $\tau a_i$. For example, when $\tau=2$ and $A=(0,3,6,\ldots)$,
the observations are made at times $0,6,12,\ldots$, but the denominator
for $n$ observations is $2n$. This convention measures the diversity
per observed instruction, with the block length accounted for.

The two optimized quantities are
\begin{equation}\label{cc-eq:LM}
 L_\tau=\sup_{A\in\mathscr S}\inf_{s\in\Ssafe_\tau}h_A(s),
 \qquad M_\tau=\inf_{s\in\Ssafe_\tau}h^*(s).
\end{equation}
They satisfy $L_\tau\le M_\tau$. If $r$ safety domains cover
$Q$, selecting the first legal block in that ordered subcover is a
Borel feedback using at most $r$ labels. It has at most $r^n$
patterns on every $n$-element window. Hence
\[
 L_\tau\le M_\tau\le\frac{\log r}{\tau}.
\]
Our main results give conditions for this elementary upper bound to
be sharp. The infimum defining $M_\tau$ is always attained:
maximal-pattern rates have the form $\log m/\tau$ for integers
$m$ (Appendix~\ref{app:minimum}).

In the three-symbol system from the introduction, the legal domains
are $K_u=\{x:x_0\ne u\}$ and $T_s=\sigma$ for every feedback.
Illegal controls may be sent to an isolated fixed point $\dagger$
outside $Q$; the resulting maps on $Q\sqcup\{\dagger\}$ are
continuous. Two domains cover $Q$. The rule $s_0(x)=1$ on
$[0]$ and $s_0(x)=0$ elsewhere has $\Patt_{s_0}(F)=2^{|F|}$:
prescribe $x_j=0$ or $1$ to obtain, respectively, label $1$ or $0$
at each $j\in F$. This verifies the upper value $\log2$ directly.
The harder question is a lower bound valid for every Borel choice.
The next section proves it on every sequence with diverging gaps.
Here $[v]$ denotes the cylinder of sequences beginning with a word $v$;
we use this notation in all subsequent shift examples.

\subsection{Observing a partition instead of a control}
A finite Borel partition may split the domain assigned to one block
into several observed atoms. Assigning a safe block to each atom gives
an invariant cover in the sense of \cite[Definition~2.8]{Kawan2013},
with the underlying cover required to be a partition. We call this an
\emph{invariant partition}. Its state map still applies the assigned
block, but its atom names may distinguish trajectories that have the
same control name. The following merging argument shows why this
additional distinction does not change the infima.

\begin{lemma}\label{cc-lem:merge}
The infima in \eqref{cc-eq:LM}, and the inner infimum for each fixed
$A$, are unchanged if finite Borel invariant partitions are allowed
and atom names are counted.
\end{lemma}
\begin{proof}
Merge the atoms assigned the same block. Their union is Borel and
lies in that block's safety domain, so the merged partition defines
the same state map. Projection of atom labels to block labels is onto
every realized finite pattern set. It can only decrease counts.
Conversely, the nonempty fibers of a selector are a finite invariant
partition whose atom and block counts agree. Taking the two infima
proves the assertion.
\end{proof}

\begin{example}[Observation labels and physical controls]\label{ex:physical-merge}
Let $Q=\{0,1\}^{\Nzero}\cup\{2^\infty\}$. Control $u=x_0$
sends $x$ to $\sigma x$ and every other control sends it to an
isolated cemetery point. Safety forces $s(x)=x_0$. The $n$-position
pattern count is $2^n+1$ for $n\ge1$. Renaming the observed label $2$ as $0$
reduces it to $2^n$, with the same exponential rate. Yet no one
physical control is safe on both $[0]$ and $\{2^\infty\}$.
Such a renaming is an observation factor, not a merged safe feedback.
Lemma~\ref{cc-lem:merge} merges only atoms assigned the same block.
\end{example}


\subsection{Finite-horizon covering}
For later comparison, let $r_{\rm inv}(N,Q)$ be the least size of a
finite set $S\subset U^N$ such that from every $x\in Q$ at least
one word in $S$ stays safe through time $N$. Set this number to
$+\infty$ if no finite spanning family exists. Invariance entropy is
\begin{equation}\label{eq:hinv-definition}
 \hinv(Q)=\limsup_{N\to\infty}\frac{\log r_{\rm inv}(N,Q)}{N}.
\end{equation}
This is the discrete-time definition in
\cite[Definition~2.2]{Kawan2013}, with natural logarithms and the
full control space $U^{\Nzero}$. The words actually produced by a
feedback form a spanning family. Conversely, a spanning family by
itself prescribes no consistent choice at later states. We will use
this distinction both in computing examples and in comparing the
two entropy notions in Section~\ref{sec:comparison}.

\section{Common successors and a closed-cover formula}\label{sec:closed-cover}

Suppose that every legal control block sends a state to the same
successor $T(x)$. A feedback then chooses labels along a fixed orbit,
just as in the excluded-letter example. Its label fibers form a Borel
partition refining the safety cover. We first solve this partition
problem on a mixing shift of finite type. The control conclusion will
follow by applying the formula to the block safety domains.

We use the counting convention of Section~\ref{sec:control} for a
finite Borel partition $\alpha$ and a Borel map $T:Q\to Q$ on a
nonempty compact metric space. In this notation
\[
 \Patt_\alpha(F)=\#\bigvee_{j\in F}T^{-j}\alpha
\]
counts nonempty atoms, or equivalently the realized $\alpha$-names
on $F$. In the present partition problem the scale is one:
\begin{align}
 h_A(T,\alpha)&=\limsup_n\frac1n
       \log\Patt_\alpha(\{a_1,\ldots,a_n\}),\label{cc-eq:seq}\\
 p_\alpha^*(n)&=\max_{|F|=n}\Patt_\alpha(F),
 &h^*(T,\alpha)&=\lim_n\frac1n\log p_\alpha^*(n).
 \label{cc-eq:pat}
\end{align}
As before, submultiplicativity gives the last limit and
$h_A(T,\alpha)\le h^*(T,\alpha)$. We restore the factor
$1/\tau$ when applying the result to control blocks.

Let $\K=(K_w)_{w\in W}$ be a finite closed cover of $Q$, and write
\[
 r=N(\K)=\min\{|I|:I\subset W,\ \bigcup_{w\in I}K_w=Q\}.
\]
The set $\mathcal B(\K)$ consists of all finite Borel partitions refining
$\K$: each atom must be contained in at least one $K_w$. Equivalently,
one may first consider all Borel maps $s:Q\to W$ with $x\in K_{s(x)}$,
and then allow arbitrary finite Borel refinements of their fibers.

Let $H$ be a primitive finite $0$--$1$ matrix, with vertex set $V$, and
let
\[
 Y_H=\{y\in V^{\Nzero}:H_{y_jy_{j+1}}=1\text{ for all }j\ge0\}.
\]
Primitivity means that all sufficiently large powers of $H$ have strictly
positive entries. Write $\sigma$ for the one-sided shift.

\begin{theorem}[Closed covers of mixing shifts]\label{cc-thm:cover}
Suppose $T:Q\to Q$ is topologically conjugate to $\sigma^p:Y_H\to Y_H$
for some primitive $H$ and integer $p\ge1$. Let $\K$ be any finite closed
cover of $Q$. For every strictly increasing $A=(a_i)$ satisfying
\begin{equation}\label{cc-eq:gaps}
 a_{i+1}-a_i\longrightarrow\infty,
\end{equation}
one has
\begin{equation}\label{cc-eq:coverformula}
 \min_{\alpha\in\mathcal B(\K)} h_A(T,\alpha)
 =\min_{\alpha\in\mathcal B(\K)}\hpat(T,\alpha)
 =\log r.
\end{equation}
An ordered difference partition of any minimum subcover attains both
minima, simultaneously for every $A$ satisfying \eqref{cc-eq:gaps}.
Consequently,
\begin{equation}\label{cc-eq:coverminimax}
 \max_A\inf_{\alpha\in\mathcal B(\K)}h_A(T,\alpha)=\log r,
\end{equation}
where the maximum ranges over all strictly increasing infinite sequences
and is attained, for example, by $a_i=(i-1)^2$.
\end{theorem}

The lower bound has two ingredients. Closedness of the cover forces
at least $r$ atoms of any refining partition to be nonmeager. Mixing
then lets us prescribe any string of these atoms at sufficiently
separated times. The separation threshold depends on the partition,
but every sequence with diverging gaps eventually exceeds it. This is
why one $A$ can be fixed before choosing the partition.

\subsection{A category and concatenation lemma}

All category statements are relative to $Q$. A nonmeager Borel set is
comeager in some nonempty open set, since Borel sets have the Baire
property. We will replace each of several nonmeager atoms by a cylinder
on which it is comeager, join the cylinders, and then remove the bad
sets. For that last step we need openness: if $F$ is continuous and
open, the preimage of a closed nowhere dense set has empty interior.
Consequently preimages of meager sets are meager.

\begin{lemma}\label{cc-lem:glue}
Under the dynamical assumptions of Theorem~\ref{cc-thm:cover}, $T$ is
continuous, open, and onto. Given finitely many nonempty cylinder sets
$C_1,\ldots,C_q$ in symbolic coordinates, there is an integer $b\ge1$
such that, for every $k\ge1$, every $t_1,\ldots,t_k\in\{1,\ldots,q\}$,
and all $0\le n_1<\cdots<n_k$ with $n_{j+1}-n_j\ge b$,
\[
 \bigcap_{j=1}^k T^{-n_j}C_{t_j}
\]
is nonempty and open.
\end{lemma}
\begin{proof}
Work on $Y_H$ through the conjugacy. The image under $\sigma$ of a
prefix cylinder of length at least two is the cylinder obtained by
removing its first symbol. The image of a one-symbol cylinder is the
union of the permitted successor cylinders. Hence $\sigma$ is open.
Every vertex has a predecessor, so every infinite path has a predecessor;
thus $\sigma$ is onto. These properties pass to powers and conjugates.

Choose within each $C_j$ a nonempty prefix cylinder, and extend their
defining words, if necessary, to a common length $\ell$. Choose $N$ with
$H^n>0$ for every $n\ge N$, and then choose $b$ with
$pb-\ell+1\ge N$. Place the word for $C_{t_j}$ at coordinate $pn_j$.
Between the end of one placed word and the start of the next there are
$p(n_{j+1}-n_j)-\ell+1$ edges to fill. Primitivity supplies a path of
that length between any two endpoint vertices. Predecessors fill the
finite segment before the first placed word, and successors provide an
infinite continuation after the last. This constructs an actual point
in the intersection. The intersection is open because it is a finite
intersection of preimages of open sets under continuous maps.
\end{proof}

\begin{lemma}\label{cc-lem:borel}
Under the same assumptions, let $\alpha$ be a finite Borel partition
with $q$ nonmeager atoms. For every $A$ satisfying \eqref{cc-eq:gaps} there
exists $m\ge1$ such that
\begin{equation}\label{cc-eq:categorycount}
 \Patt_\alpha(\{a_1,\ldots,a_n\})\ge q^{n-m+1}\qquad(n\ge m).
\end{equation}
In particular $h_A(T,\alpha)\ge\log q$.
\end{lemma}
\begin{proof}
We first choose cylinders in which the selected atoms occupy a
comeager set. We then join those cylinders along the prescribed times;
the discarded subsets remain meager under pullback.

List the nonmeager atoms as $D_1,\ldots,D_q$. The Baire property
provides a nonempty cylinder $C_j$ for which $C_j\setminus D_j$ is
meager. Let $b$ be supplied by Lemma~\ref{cc-lem:glue}, and choose $m$ so
that $a_{i+1}-a_i\ge b$ for all $i\ge m$.

Fix $n\ge m$ and any target labels $t_m,\ldots,t_n$. The set
\[
 G=\bigcap_{i=m}^n T^{-a_i}C_{t_i}
\]
is nonempty and open. Each power $T^{a_i}$ is continuous and open, so
\[
 \bigcup_{i=m}^n T^{-a_i}(C_{t_i}\setminus D_{t_i})
\]
is meager and cannot cover $G$. A point outside this union realizes
all the specified atoms at the sampled times. Distinct target label
strings give distinct patterns; restriction from the full window to
this tail is onto. This proves \eqref{cc-eq:categorycount}. Dividing its
logarithm by $n$ and taking a limit inferior proves the asserted bound.
\end{proof}

\subsection{Proof of the exact formula}

\begin{proof}[Proof of Theorem~\ref{cc-thm:cover}]
Let $\alpha\in\mathcal B(\K)$. Assign to each atom $D$ one index
$w(D)$ with $D\subset K_{w(D)}$. Let $J$ be the set of indices assigned
to the nonmeager atoms. If the closed sets $(K_w)_{w\in J}$ did not
cover $Q$, their complement would be a nonempty open set covered by
the finitely many meager atoms of $\alpha$, which is impossible.
Thus $|J|\ge r$, and the number $q$ of nonmeager atoms is at least $r$.
Lemma~\ref{cc-lem:borel} gives
\[
 h_A(T,\alpha)\ge\log q\ge\log r
\]
for every $A$ satisfying \eqref{cc-eq:gaps}. Closedness of the cover
members is used precisely in this passage from category to a full cover.

Choose a minimum subcover $K_{w_1},\ldots,K_{w_r}$, and set
\begin{equation}\label{cc-eq:ordered}
 E_j=K_{w_j}\setminus\bigcup_{k<j}K_{w_k},\qquad
 \alpha_0=\{E_1,\ldots,E_r\}.
\end{equation}
Every $E_j$ is Borel and is nonempty, since a minimum subcover has no
redundant member. The partition refines $\K$ and has $r$ atoms. Hence
$\Patt_{\alpha_0}(F)\le r^{|F|}$ for all finite $F$, giving
$h_A(T,\alpha_0)\le\hpat(T,\alpha_0)\le\log r$ for every $A$.
The established lower bound and $\hpat\ge h_A$ yield
\eqref{cc-eq:coverformula} and simultaneous attainment.
For arbitrary $A$, the partition $\alpha_0$ bounds the inner infimum
in \eqref{cc-eq:coverminimax} by $\log r$. Any $A$ satisfying
\eqref{cc-eq:gaps} gives the reverse bound and attains the maximum.
\end{proof}


\subsection{An excluded-letter cover}

For any finite cover $\K$ define its combinatorial sequence entropy by
\[
 h_A(T,\K)=\limsup_{n\to\infty}\frac1n
 \log N\left(\bigvee_{i=1}^n T^{-a_i}\K\right).
\]
For a partition this is \eqref{cc-eq:seq}. For overlapping sets, covering
each finite orbit window need not correspond to one fixed refining
partition used at every time.

\begin{corollary}\label{cc-cor:covergap}
Let $q\ge3$, $Q=\{0,\ldots,q-1\}^{\Nzero}$, $T=\sigma$, and
\[
 K_j=\{x:x_0\ne j\}\quad(0\le j<q),\qquad \K=(K_j)_{j=0}^{q-1}.
\]
For every increasing infinite $A$,
\begin{equation}\label{cc-eq:covervalue}
 h_A(\sigma,\K)=\log\frac q{q-1}.
\end{equation}
For every such $A$ with diverging successive gaps,
\begin{equation}\label{cc-eq:refvalue}
 \min_{\alpha\in\mathcal B(\K)}h_A(\sigma,\alpha)=\log2
 >h_A(\sigma,\K).
\end{equation}
\end{corollary}
\begin{proof}
The cover has no single member equal to $Q$, and any two distinct
members cover $Q$. Thus $r=2$, and \eqref{cc-eq:refvalue} follows from
Theorem~\ref{cc-thm:cover} once \eqref{cc-eq:covervalue} is proved.

For a window $F$ of size $n$, the full shift has all $q^n$ possible
input vectors. Each member of the joined cover specifies one excluded
letter at each coordinate, and covers $(q-1)^n$ of these vectors.
Consequently the minimum number $c(n)$ of such boxes is at least
$(q/(q-1))^n$. For completeness, the following elementary probabilistic
covering bound gives the matching exponent. Choose $R$ excluded-letter
vectors independently and uniformly from the $q^n$ possibilities. A
fixed input vector is left uncovered with probability at most
\[
 \left(1-\left(\frac{q-1}{q}\right)^n\right)^R
 \le \exp\left(-R\left(\frac{q-1}{q}\right)^n\right).
\]
Taking
$R=\left\lceil(q/(q-1))^n(n\log q+1)\right\rceil$ makes the sum of
these probabilities over all input vectors smaller than one. Some
choice therefore covers every vector. We have proved
\[
 \left(\frac q{q-1}\right)^n\le c(n)
 \le\left\lceil\left(\frac q{q-1}\right)^n(n\log q+1)\right\rceil.
\]
The count depends only on $n$, not on the positions of the sampled
coordinates, so its exponential rate proves \eqref{cc-eq:covervalue}.
\end{proof}

The excluded-letter cover and its covering exponent are already present
in Romagnoli \cite{Romagnoli2007}; the calculation above is included to
identify the sparse-time quantity without an external dependency.
The consequence of Theorem~\ref{cc-thm:cover} is the exact minimum over
\emph{all Borel} refinements for a sequence fixed in advance. It does not
assert that every Borel refinement has consecutive-time entropy at least
$\log2$, or that every sampling sequence gives that lower bound.


\subsection{Application to safe feedbacks}\label{sec:common-successor}

For a fixed period, the common-successor hypothesis identifies every
feedback orbit with the same dynamical system. The preceding cover
formula therefore applies to the control-label fibers.

\begin{corollary}[Safe feedbacks with a common successor]\label{cc-thm:control}
In the setting of Section~\ref{sec:selectors}, assume that there is
a map $T:Q\to Q$ such that
\begin{equation}\label{cc-eq:common}
 g_w(x)=T(x)\qquad\text{for every }w\in W\text{ and every }x\in K_w,
\end{equation}
and that $T$ is conjugate to $\sigma^p$ on a mixing one-sided shift of
finite type. With $r=N((K_w)_{w\in W})$, for every $A$ satisfying
\eqref{cc-eq:gaps} one has
\begin{equation}\label{cc-eq:controlformula}
 \min_{s\in\Ssafe_\tau}h_A(s)
 =\min_{s\in\Ssafe_\tau}\hpat(s)
 =L_\tau=M_\tau=\frac{\log r}{\tau}.
\end{equation}
The maximum defining $L_\tau$ is attained by every such $A$, and an
ordered minimum-subcover selector attains all the indicated minima.
The same minimum and minimax values hold when all finite Borel safe
partitions, including every finite Borel refinement, are allowed. Every
refinement satisfies the stated lower bound; no upper bound for each
individual refinement is asserted.
\end{corollary}
\begin{proof}
For every legal selector, \eqref{cc-eq:common} gives $T_s=T$ pointwise.
Its fibers form a finite Borel refinement of $(K_w)_{w\in W}$.
Theorem~\ref{cc-thm:cover}, with normalization by $\tau$, gives the lower
bound at every $A$ with diverging gaps. If $w_1,\ldots,w_r$ form a
minimum subcover, define $s_0(x)=w_j$ on the sets $E_j$ of
\eqref{cc-eq:ordered}. This is Borel and legal at every point. Each selected
block stays in $Q$ at its intermediate times by \eqref{cc-eq:safe}, and
the endpoint again lies in $Q$, so recursion gives an infinite safe
closed-loop orbit for every initial state. The selector uses only $r$
blocks and thus gives the upper bound for every window and every
sampling sequence. This proves attainment, \eqref{cc-eq:controlformula},
and the outer maximum. Lemma~\ref{cc-lem:merge} proves the last assertion.
\end{proof}

At period one this gives $L_1=M_1=\log2$ for the excluded-letter
system from the introduction. At longer periods, the cover cardinality
must be computed for whole blocks, as the following example shows.

\subsubsection{An overlapping example at period two}

Take
\[
 Q=\{0,1,2\}^{\Nzero},\quad X=Q\sqcup\{\dagger\},\quad
 U=\{0,1,2\},\quad\tau=2,
\]
where the isolated point $\dagger$ lies outside $Q$. Define
\[
 f_u(x)=
 \begin{cases}\sigma x,&x\in Q,\ x_0\ne u,\\
 \dagger,&x\in Q,\ x_0=u,\end{cases}
 \qquad f_u(\dagger)=\dagger.
\]
Each $f_u$ is continuous since the cases are clopen. For $w=(a,b)$,
\[
 K_{(a,b)}=\{x:x_0\ne a,\ x_1\ne b\},\qquad
 g_{(a,b)}|_{K_{(a,b)}}=\sigma^2.
\]
All nine control blocks are available to Borel feedbacks, and each
state permits four of them. A safety domain covers four of the nine
two-symbol input blocks, so two domains cannot cover $Q$. The three
diagonal blocks $(0,0),(1,1),(2,2)$ do cover $Q$, since some symbol
differs from both $x_0$ and $x_1$. Hence $r=3$, and
\[
 L_2=M_2=\frac{\log3}{2}.
\]
An attaining selector is
\[
 c(x)=\min(\{0,1,2\}\setminus\{x_0,x_1\}),\qquad
 s_0(x)=(c(x),c(x)).
\]
It is clopen and legal everywhere. Its actual block-boundary states
are $\sigma^{2k}x$. To prescribe any sequence of diagonal labels, use
successive input blocks $11$, $00$, and $01$ for the desired labels
$0$, $1$, and $2$, respectively. Therefore its actual name set is the
full shift on these three blocks. The lower bound, however, holds for
every legal Borel rule using any of the nine blocks.

The common object here is the state orbit; there need not be a
nonconstant label observation shared by all feedbacks. To see this,
suppose a map
$\rho:U^2\to V'$ and an observation $F:Q\to V'$ satisfy
$\rho(s(x))=F(x)$ for every legal Borel selector and every $x$.
Any two blocks $(a,b)$ and $(a',b')$ are both legal at a state beginning
with $ij$, where $i\notin\{a,a'\}$ and $j\notin\{b,b'\}$.
A legal baseline selector may be changed at this single state to use
either block, while remaining Borel and legal. Hence their $\rho$
values agree. Every pair of blocks is legal together somewhere, so $\rho$
is constant. Thus the argument uses refinements of an
overlapping cover, rather than a label forced at each state.


\subsection{Why the successor map matters}\label{sec:different}
Changing the successor can prevent a region from being visited again.
The minimum safety cover still counts that region, but its label may
contribute only to the initial part of a name. The next example has
$r=3$ and $L_1=M_1=\log2$. Thus continuity and openness of the
individual branches do not replace the common-successor assumption.

Let $Q=\{0,1\}^{\Nzero}$, let $\mathcal C=\{0,10,11\}$ be a set of
finite words, and take $U=\mathcal C\times\{0,1\}$ and $\tau=1$.
For $e\in\{0,1\}$ let $P_e$ add $e$ modulo two to every coordinate.
On $X=Q\sqcup\{\dagger\}$ define
\begin{equation}\label{cc-eq:boundarysystem}
 f_{(c,e)}(x)=
 \begin{cases}P_e\sigma x,&x\in[c],\\
 \dagger,&x\in Q\setminus[c],\end{cases}
 \qquad f_{(c,e)}(\dagger)=\dagger.
\end{equation}
The cylinders $[0],[10],[11]$ are disjoint and cover $Q$. Each state
has exactly two legal controls with different successors. The global
maps $f_{(c,e)}$ are continuous, and their legal restrictions are open.
The safety domains are $K_{(c,e)}=[c]$, and $r=3$.

\begin{proposition}\label{cc-prop:boundary}
For the system \eqref{cc-eq:boundarysystem}, over all legal Borel feedbacks
and arbitrary Borel mixtures,
\begin{equation}\label{cc-eq:boundaryvalue}
 \min_s h_{A_0}(s)=\min_s\hpat(s)=
 \max_A\inf_s h_A(s)=\log2<\log r,
 \qquad A_0=(0,1,2,\ldots).
\end{equation}
There is one clopen legal feedback $s_{\mathrm f}$ satisfying
$h_A(s_{\mathrm f})=\log2$ for every increasing infinite $A$.
The same full-class values hold if arbitrary finite Borel atom
refinements are allowed.
\end{proposition}
\begin{proof}
For the upper bound we choose the flip that makes the next first
coordinate zero. After that first step only two labels remain. For
the lower bound we recover the input prefix from any actual control
word, whatever feedback produced it.

Let $c(x)$ be the unique member of $\mathcal C$ with $x\in[c(x)]$,
and set
\[
 s_{\mathrm f}(x)=(c(x),x_1),\qquad
 T_{\mathrm f}x=P_{x_1}\sigma x.
\]
This is a clopen selector, legal at every state. Its four nonempty
initial control fibers, for labels $(0,0),(0,1),(10,0),(11,1)$, are
$[00],[01],[10],[11]$, respectively. For all $n\ge1$, direct induction
gives
\begin{equation}\label{cc-eq:fliporbit}
 T_{\mathrm f}^n x=P_{x_n}\sigma^n x,
 \qquad s_{\mathrm f}(T_{\mathrm f}^n x)
        =(0,x_{n+1}\mathbin\oplus x_n),
\end{equation}
where $\oplus$ is addition modulo two. In the induction step the
chosen flip is $x_{n+1}\oplus x_n$; it cancels the previous flip and
leaves $P_{x_{n+1}}\sigma^{n+1}x$. The first coordinate is always zero
after time zero, which also proves all-time safety.

For every finite nonempty window $F$, with $n=|F|$, one has
\begin{equation}\label{cc-eq:flipcount}
 \Patt_{s_{\mathrm f}}(F)=
 \begin{cases}2^n,&0\notin F,\\2^{n+1},&0\in F.\end{cases}
\end{equation}
The four initial labels and two later labels give the upper bounds.
To prove the lower bounds, prescribe any of the four initial labels
if $0\in F$, thereby fixing $(x_0,x_1)$; if $0\notin F$, set
$x_0=x_1=0$. At positive sampled times prescribe the desired second
components $d_k$, and put $d_k=0$ at the other positive times.
Recursively set $x_{k+1}=x_k\oplus d_k$ for every $k\ge1$.
This constructs a full state in $Q$ whose actual trajectory under the
same selector has the required sampled controls by
\eqref{cc-eq:fliporbit}. This realizes every prescribed pattern on
an infinite closed-loop trajectory. Equation~\eqref{cc-eq:flipcount}
implies $h_A(s_{\mathrm f})=\hpat(s_{\mathrm f})=\log2$ for every $A$.

For the opposite bound fix any legal Borel selector $s$. If its first
$n$ actual controls from $x$ are $(c_k,e_k)$, let
$E_0=0$ and $E_k=e_0\oplus\cdots\oplus e_{k-1}$. Since shifts commute
with flips,
\[
 T_s^k x=P_{E_k}\sigma^k x.
\]
Let $b(c)$ be the first bit of the word $c$. Safety at time $k$ gives
\begin{equation}\label{cc-eq:decode}
 x_k=b(c_k)\oplus E_k\qquad(0\le k<n).
\end{equation}
Thus one actual length-$n$ control word determines at most one input
prefix of length $n$. Every binary input prefix has an extension in
$Q$, and the feedback is defined and safe on every such extension.
Distinct prefixes must yield distinct actual control words. Hence
\[
 \Patt_s(\{0,\ldots,n-1\})\ge2^n
\]
for every $n$ and every legal $s$. This proves $h_{A_0}(s)\ge\log2$
and $\hpat(s)\ge\log2$. Combining these inequalities with
$s_{\mathrm f}$ proves \eqref{cc-eq:boundaryvalue}, including attainment.
Lemma~\ref{cc-lem:merge} handles atom refinements.
\end{proof}

The obstruction has a precise orbit interpretation. In this example
$T_{\mathrm f}(Q)=[0]$: for the reverse inclusion, given $y\in[0]$,
take $x=0y$. Therefore neither $[10]$ nor $[11]$ is visited at a
positive time. The initial fibers with those labels are nonmeager,
but their labels cannot be freely concatenated in the future.
The map $T_{\mathrm f}$ is itself continuous and open, as follows by
restricting to the four clopen two-symbol cylinders. It is not onto
$Q$.

Even preservation of meager preimages does not repair this problem.
For any selector $s$ and any $n$, the value of $T_s^n x$ is either
$\sigma^n x$ or $P_1\sigma^n x$. Hence, for a meager set $N$,
\[
 (T_s^n)^{-1}N\subset
 (\sigma^n)^{-1}N\ \cup\ (P_1\sigma^n)^{-1}N
\]
is meager. What fails is the nonempty open intersection used in
Lemma~\ref{cc-lem:glue} for the \emph{selected} dynamics. Also, the
same $s_{\mathrm f}$ gives
\[
 \inf_s h_A(s)\le\log2=\log r-\log(3/2)
 \quad\text{for every }A.
\]
Thus the deficit below $\log r$ persists for every sampling sequence,
while $L_1=M_1$ still holds. The exact inner minimum at other fixed
sequences is not determined by this calculation.




\section{Inverse capacity when the feedback changes the orbit}\label{sec:capacity}

The preceding example shows why the safety cover alone is insufficient
when legal controls have different successors. We now use a reference
measure to bound the size of the initial-state sets carrying control
names. The measure is only a tool for counting; it need not be
invariant under any feedback.

The estimate concerns sets of initial states whose controls belong to
prescribed lists of $k-1$ labels. If each such set has exponentially
small mass, exponentially many lists are needed to cover the names of
any feedback. A finite covering lemma then supplies a set of positions
on which all choices of $k$ labels occur. The mass bound is uniform
over feedbacks, but these positions can depend on the rule. This is
why the argument determines $M_\tau$ without yet determining $L_\tau$.

\subsection{Safe inverse-transfer operators}
Fix a finite block set $W$ of size $q$ at period $\tau$, with
closed safety domains $K_u$ covering $Q$ and continuous endpoint
maps $g_u:K_u\to Q$, as in \eqref{cc-eq:safe}. All optimizations
in this section use the Borel selectors with values in this $W$.
Let $\mu$ be a full-support Borel probability on $Q$ and assume
\begin{equation}\label{eq:cap-density}
 \nu_u=(g_u)_*(\mu|_{K_u})\ll\mu,\qquad
 \rho_u=\frac{d\nu_u}{d\mu}\in L^\infty(\mu).
\end{equation}
The function $\rho_u$ measures how much source mass the safe branch
can send to a target region. To follow this mass through several
branches, define $L_u$ by
\begin{equation}\label{eq:cap-transfer}
 \int hL_uf\,d\mu=\int_{K_u}h(g_ux)f(x)\,d\mu(x)
 \quad(h,f\text{ bounded Borel}).
\end{equation}
Thus $L_uf$ is the density after pushing $f\mu$ through branch $u$
and restricting the source to its safety domain. The signed
pushforward is dominated in total variation by $\|f\|_\infty\nu_u$;
its density exists and satisfies
$|L_uf|\le\|f\|_\infty\rho_u$. Changing $f$ on a null set
does not change the pushforward, so $L_u$ is a positive bounded
operator on $L^\infty(\mu)$ and $L_u\one=\rho_u$.

Fix $2\le k\le q$ and set $l=k-1$. A list box of length $n$
is a product $S_0\times\cdots\times S_{n-1}$ with $S_j\subset W$
and $|S_j|=l$: it allows any label in $S_j$ at time $j$.
For one list $S$, let $A_S=\sum_{u\in S}L_u$. For example, two
successive lists $S_0,S_1$ give
\[
 A_{S_1}A_{S_0}\one
   =\sum_{u\in S_0,\ v\in S_1}L_vL_u\one.
\]
Each summand is the terminal density obtained by following one
safe two-block path. A feedback uses only part of the source set
available to each path. Pushing forward the initial states whose
first two labels belong to these lists therefore gives a terminal
density bounded by the displayed sum. Integrating gives an upper
bound on their initial
mass; since $\mu$ is a probability, the $L^\infty$ norm bounds that
integral as well. We use this norm to estimate list boxes of any length:

\begin{equation}\label{eq:cap-bn}
 b_0=1,\qquad
 b_n=\max_{|S_0|=\cdots=|S_{n-1}|=l}
       \|A_{S_{n-1}}\cdots A_{S_0}\one\|_\infty.
\end{equation}
The first list acts on the right. A history in \eqref{eq:cap-bn}
chooses one list at each time for the whole state space; it does not
choose a different list at each target point. This distinction will
matter in Section~\ref{subsec:product-example}.

For a positive operator $P$, $\|P\|_{\infty\to\infty}
=\|P\one\|_\infty$. Hence $b_{n+m}\le b_nb_m$. In particular,
one strict bound $b_h<1$ controls all lengths: choose
$b_h<\theta^h<1$ and put
\begin{equation}\label{eq:cap-exp}
 H=\max_{0\le j<h}b_j\theta^{-j},\qquad
 b_n\le H\theta^n\quad(n\ge0).
\end{equation}
Indeed, write $n=ah+j$ and use submultiplicativity. The integer
$h$ is a length used to verify the estimate; the control period
remains $\tau$.

A set of words $G\subset W^n$ has a \emph{$k$-valued cube projection}
on $J\subset\{0,\ldots,n-1\}$ if, for some $V\subset W$ with
$|V|=k$, its projection onto $J$ contains $V^J$. Thus every assignment
of those $k$ labels to the positions in $J$ occurs in a word of $G$.
The theorem produces such projections inside the actual feedback
language, with a number of positions proportional to $n$.


\begin{theorem}[Product inverse-capacity criterion]\label{thm:capacity-main}
Assume \eqref{eq:cap-density}. If $b_h<1$ for some $h\ge1$, then
there are $\alpha>0$ and $N$, depending only on the branches,
$\mu,k$, such that, for every legal Borel feedback $s$ and every
$n\ge N$, its actual length-$n$ names have a $k$-valued cube
projection on at least $\alpha n$ coordinates. Consequently
\begin{equation}\label{eq:cap-pattern}
 p_s^*(m)\ge k^m\quad(m\ge1),\qquad
 h^*(s)\ge\frac{\log k}{\tau}.
\end{equation}
The same lower bound holds for all finite Borel atom refinements.
If $k=r\ge2$ is the minimum number of safety domains covering $Q$,
then
\begin{equation}\label{eq:cap-M}
 M_\tau=\min_s h^*(s)=\frac{\log r}{\tau},
\end{equation}
both over unrefined selectors and over all finite Borel safe partitions.
\end{theorem}

\subsection{From actual fibers to list covers}
For a fixed feedback, group initial states by their first $n$ controls:
let
\[
 E_w(s)=\{x:s(T_s^jx)=w_j\ (0\le j<n)\},\qquad
 F_n(s)=\{w\in W^n:E_w(s)\ne\varnothing\}.
\]
The sets $E_w(s)$ are the actual fibers. They are Borel, pairwise
disjoint, and cover $Q$. A nonempty fiber is included even when it
has measure zero.
For $w=w_0\cdots w_{n-1}$ put
$L_w=L_{w_{n-1}}\cdots L_{w_0}$ and
$P_F=\sum_{w\in F}L_w$.

\begin{lemma}[Actual-fiber capacity]\label{lem:fiber-cap}
For every legal Borel feedback and every $n\ge1$,
\begin{equation}\label{eq:fiber-mass}
 1\le\int P_{F_n(s)}\one\,d\mu.
\end{equation}
If $F_n(s)$ is covered by $t$ boxes
$S^a_0\times\cdots\times S^a_{n-1}$ whose lists have size at
most $l$, then $1\le t b_n$.
\end{lemma}
\begin{proof}
Follow the mass in one fiber through its prescribed branches. Safety
means that none is lost at a restriction to $K_{w_j}$. Dropping the
fiber restriction at the start can only increase the resulting density.

Fix a nonempty fiber $E_w$ and let $\eta_j$ be the pushforward
of $\mu|_{E_w}$ after its first $j$ actual blocks. Set
$f_0=\one_{E_w}$. We claim inductively that
$\eta_j=f_j\mu$ and $f_{j+1}=L_{w_j}f_j$. Indeed,
$\eta_j$ is concentrated on $K_{w_j}$, and for every bounded Borel $h$,
\[
 \int h\,d\eta_{j+1}
 =\int_{K_{w_j}}h(g_{w_j}x)f_j(x)\,d\mu(x)
 =\int hL_{w_j}f_j\,d\mu.
\]
The boundedness of each $L_u$ permits the induction. None of these
restrictions removes mass, because all states in $E_w$ follow the
specified safe word. Thus $\int L_w\one_{E_w}\,d\mu=\mu(E_w)$.
Positivity gives $L_w\one_{E_w}\le L_w\one$; summing over the
disjoint fibers proves \eqref{eq:fiber-mass}. The set $F_n(s)$ includes
null but nonempty fibers as well.

Complete every nonempty list to size $l$. Distributivity and positivity
give
\[
 P_{F_n(s)}\one\le\sum_{a=1}^t
 A_{S^a_{n-1}}\cdots A_{S^a_0}\one.
\]
Repeated boxes and overlaps only enlarge this upper bound. Unsafe
words in a box are subject to the original safety gates in $L_u$.
Integrating and using \eqref{eq:cap-bn} yields $1\le tb_n$.
\end{proof}

\subsection{The finite combinatorial input and its application}
The preceding mass estimate gives a lower bound on the number of
list boxes needed to cover the names. The finite covering lemma we
will use is stated for a different kind of box: at each position it
omits one letter from a fixed $k$-element alphabet. We compare these
two covering numbers before applying the lemma.

For $F\subset W^n$, let $C_l(F)$ be its minimum covering number by
$l$-list boxes. For a $k$-element set $V\subset W$, let $C_V(F)$
be the minimum number of boxes of the form
\[
 \prod_{j=0}^{n-1}(W\setminus\{v_j\}),\qquad v_j\in V,
\]
needed to cover $F$. This is a cover, since $k\ge2$.

We use the following finite covering lemma. For every
$k\ge2$ and $b>0$ there are $c>0$ and $N$ such that if
$G\subset\{0,1,\ldots,k\}^n$, $n\ge N$, needs at least
$e^{bn}$ boxes omitting one nonzero letter per coordinate, then some
$J\subset\{0,\ldots,n-1\}$ with $|J|\ge cn$ has
\begin{equation}\label{eq:HY-cube}
 G|_J\supset\{1,\ldots,k\}^{J}.
\end{equation}
This is Kerr--Li's finite covering lemma
\cite[Lemma~3.3]{KL2007}. Huang and Ye recover the same conclusion
from their sharper estimates in \cite[Lemma~4.1 and Remark~4.2,
pp.~1701--1705]{HuangYe2009}. Its constants depend only on $b,k$;
its use here requires no dynamical assumptions on the closed loop.

To pass from $C_l$ to one of the numbers $C_V$, choose a minimum
omitted-letter cover for every $V$ and intersect all these covers. The intersections cover $F$, with at most
$\prod_V C_V(F)$ members. In a nonempty intersection each coordinate
contains at most $k-1$ letters: if it contained a $k$-set $V_0$, its
factor from the $V_0$ cover would have omitted a letter of $V_0$.
Hence
\begin{equation}\label{eq:cap-intersect}
 C_l(F)\le\prod_{V\in\binom Wk}C_V(F).
\end{equation}
For actual names Lemma~\ref{lem:fiber-cap} gives
$C_l(F_n(s))\ge b_n^{-1}$. Under \eqref{eq:cap-exp}, for all
sufficiently large $n$ this lower bound is at least
$e^{\eta n/2}$, where $\eta=-\log\theta>0$. With
$J_0=\binom qk$, \eqref{eq:cap-intersect} implies that some $V$
satisfies
\[
 C_V(F_n(s))\ge e^{\eta n/(2J_0)}.
\]
Code the members of $V$ by $1,\ldots,k$ and all other letters by
$0$. Omitted-letter boxes pull back exactly to the boxes defining
$C_V$, so the coded set has precisely that covering number. Applying
\eqref{eq:HY-cube} and undoing this coding yields
$F_n(s)|_J\supset V^J$. No distinct members of $V$ were merged.

\begin{proof}[Proof of Theorem~\ref{thm:capacity-main}]
The preceding argument supplies $\alpha,N$ independent of $s$.
Given $m$, choose $n$ so large that $\alpha n\ge m$ and retain
$m$ of the independent positions. All $k^m$ patterns are actual
feedback names, proving \eqref{eq:cap-pattern} with the prescribed
normalization by $m\tau$.

For $k=r$, choose a minimum subcover $K_{u_1},\ldots,K_{u_r}$ and
select $u_j$ on $K_{u_j}\setminus\bigcup_{i<j}K_{u_i}$. This is a
Borel state function, legal at every point. Each chosen block stays
safe at its intermediate times and ends in $Q$; applying the same
rule at the next state gives infinite safety. It uses $r$ labels and
has at most $r^m$ patterns on any $m$ positions. This proves
\eqref{eq:cap-M} and attainment. For a refined partition, projection
to its physical control labels is onto the actual coarse patterns
and preserves the dynamics. The lower bound therefore passes to
every refinement, while the unrefined selector remains an available
upper-bound witness in the enlarged class.
\end{proof}

The conclusion controls $M_\tau$. To control $L_\tau$ by this
method, one would also need to choose the independent positions in
common for all feedbacks. The covering lemma provides no such choice:
its $J$ and $V$ may depend on both $s$ and $n$.

\subsection{Single-step scalar and joint-density criteria}
Define
\begin{equation}\label{eq:cap-Theta}
 \Theta_k=\left\|\max_{|S|=k-1}\sum_{u\in S}\rho_u\right\|_\infty.
\end{equation}
If $\Theta_k<1$, then $b_1=\Theta_k<1$ and
Theorem~\ref{thm:capacity-main} applies. If the branch capacities obey
$\mu(K_u\cap g_u^{-1}E)\le c_u\mu(E)$ for every Borel $E$,
then $\rho_u\le c_u$ almost everywhere, and the sum of the largest
$k-1$ of the $c_u$ bounds $\Theta_k$. Thus a scalar-capacity sum
smaller than one is a sufficient special case. The order of taking
the branch supremum and summing matters: joint densities may pass
the test while the individual optimal scalar bounds do not.

Appendix~\ref{app:weighted} gives a direct quantitative estimate in
this single-step case. The strict inequality is essential:
Proposition~\ref{prop:critical} gives systems with $r=3$,
$\Theta_3\ge1$ arbitrarily close to one, and $M_1=\log2$.

\subsection{The excluded-letter system with different successors}
Return to the excluded-letter cover, but add the chosen control to
every coordinate of the shifted state. The safety cover is unchanged;
the trajectory now depends on the feedback. This gives a simple
application of the capacity theorem and shows what information it
leaves open.

Let $q\ge3$, $G=\mathbb Z/q\mathbb Z$, and $Q=G^{\Nzero}$.
Write $P_u$ for coordinatewise addition of $u$. On
$X=Q\sqcup\{\dagger\}$ use the continuous controls
\[
 f_u(x)=P_u\sigma x\quad\text{if }x_0\ne u,
 \qquad f_u(x)=\dagger\quad\text{otherwise},
 \qquad f_u(\dagger)=\dagger.
\]
Their safety domains are the excluded-letter cylinders
$K_u=\{x:x_0\ne u\}$, so the minimum cover cardinality is two.

\begin{corollary}\label{cor:cyclic}
For this system, at period one and over all legal Borel feedbacks
and finite atom refinements,
\[
 M_1=\min_s h^*(s)=\log2,\qquad
 \log\frac q{q-1}\le L_1\le\log2,\qquad
 \hinv(Q)=\log\frac q{q-1}.
\]
There is a clopen feedback with $\Patt_s(F)=2^{|F|}$ for every
finite window $F$.
\end{corollary}
\begin{proof}
Use the uniform Bernoulli probability on $Q$. The safe inverse
operator is
\[
 (L_uf)(y)=\frac1q\sum_{a\ne u}f(aP_{-u}y),
 \qquad L_u\one=\frac{q-1}{q}\one.
\]
Products of any single-label lists therefore have norm
$((q-1)/q)^n$. Theorem~\ref{thm:capacity-main}, with $k=r=2$,
gives $M_1=\log2$.

For a control word $w=(w_0,\ldots,w_{n-1})$ produced by $s$, put
$E_k=\sum_{j<k}w_j$ and $v_k=w_k-E_k$, with arithmetic in $G$.
The actual orbit satisfies $T_s^kx=P_{E_k}\sigma^kx$, so its
initial fiber lies in
\[
 K(w)=\{x:x_k\ne v_k\ (0\le k<n)\}.
\]
This is also the complete open-loop safety domain of $w$.
Each such domain contains exactly $(q-1)^n$ input prefixes.
Every input prefix extends to a state with an infinite safe
feedback trajectory; hence the actual length-$n$ word count is
at least $(q/(q-1))^n$. Consecutive sampling gives the stated
lower bound for $L_1$.

Define $s_0(x)=1$ if $x_0=0$, and $s_0(x)=0$ otherwise.
This is clopen and safe on all of $Q$. Given any infinite binary
sequence $d$, let $E_k=\sum_{j<k}d_j$ and let $b_k=0$ when
$d_k=1$, $b_k=1$ when $d_k=0$. The state with coordinates
$x_k=b_k-E_k$ has current first coordinate $b_k$ at time $k$,
by induction. Its feedback name is therefore $d$.
This proves the finite-window equality and the upper bound for
$L_1$ with the same feedback for every sampling sequence.

The map $w\mapsto v$ is triangular and bijective: recover
$w_k=v_k+E_k$, $E_{k+1}=2E_k+v_k$, starting with $E_0=0$.
This requires no inverse of $2$, and includes even $q$.
Thus the open-loop spanning number is the least number $c_q(n)$
of boxes $\prod_{k<n}(G\setminus\{v_k\})$ covering $G^n$.
Writing $p=((q-1)/q)^n$, the counting lower bound and the
elementary random-cover upper bound give
\[
 p^{-1}\le c_q(n)\le
 \left\lceil p^{-1}(n\log q+1)\right\rceil.
\]
Indeed, $t=\lceil p^{-1}(n\log q+1)\rceil$ independent uniform
boxes leave some input uncovered with probability at most
$q^n(1-p)^t\le q^ne^{-pt}\le e^{-1}<1$.
A covering choice exists, and removal of repetitions can only
decrease its size. Taking normalized logarithms proves the
invariance-entropy formula. Lemma~\ref{cc-lem:merge} supplies
the refinement statement.
\end{proof}

The safe spanning number is a lower bound for the actual feedback
word count, not an upper bound. For example, when $q=3$, the
three safe words $00,12,21$ cover all two-symbol inputs: their
excluded vectors are $00,11,22$. The clopen rule
$s(x)=\min(G\setminus\{x_0,x_1\})$ repeatedly selects the
first control of a covering word and is safe everywhere.
Starting from $x=(1,1,0,0,\ldots)$, however, it selects $0$
and then $2$, giving the word $02$ outside that cover.
The subsequent state is $2^\infty$ and all subsequent controls
are $0$, so the witness has an infinite safe continuation.
Its prefix fibers obey
\[
 H_s(uw)=s^{-1}(u)\cap(P_u\sigma)^{-1}H_s(w),
 \qquad H_s(w)\subset K(w).
\]
These compatibility equations are absent from a separately
chosen spanning family. The lower and upper bounds for $L_1$
remain distinct; neither their equality nor a strict minimax
gap is asserted for this cyclic system.


\subsection{A smooth joint-density example}\label{subsec:smooth}
The joint-density test can be sharper than taking a separate
supremum for each branch: peaks at different target points need not
add. On $Q=[0,1]$ take three full inverse branches with
\begin{align}
 l_u&=\frac{19u}{60},\qquad
 \rho_u(y)=\frac{11}{30}+\frac7{30}
            \cos(2\pi(y-u/3)),\label{eq:smooth-rho}\\
 \psi_u(y)&=l_u+\frac{11}{30}y+
 \frac7{60\pi}\big[\sin(2\pi(y-u/3))-\sin(-2\pi u/3)\big].
 \nonumber
\end{align}
Here $g_u=\psi_u^{-1}$, for $u=0,1,2$.
Their derivatives lie between $2/15$ and $3/5$, so the branches
are strictly increasing and expanding in the forward direction.
The closed safety intervals are
$[0,22/60]$, $[19/60,41/60]$, and $[38/60,1]$.
They cover $Q$ and each is needed: the middle interval alone covers
the open gap between $22/60$ and $38/60$, while the two endpoints
require the first and third intervals. Hence $r=3$.

With Lebesgue probability, the densities are exactly
$\psi'_u=\rho_u$. The three cosine terms sum to zero. The sum of
two is the negative of the remaining cosine, whose maximum is one;
therefore
\[
 \Theta_3=\frac{22}{30}+\frac7{30}=\frac{29}{30}<1.
\]
Each branch has optimal scalar bound $3/5$, since its continuous
density reaches that value. The largest two scalar bounds sum to
$6/5$. The joint criterion consequently gives $M_1=\log3$ although
that Lebesgue scalar criterion fails.

To place these interval branches in a continuous compact ambient
system, take $X=[-1,2]$. Extend each $g_u$ to the left of its
safety interval $[a_u,b_u]$ by $-\min\{a_u-x,1\}$ and to its
right by $1+\min\{x-b_u,1\}$. At the endpoints these formulas
match $0$ and $1$. Outside the closed safety interval the image
lies outside $Q$, so this extension changes neither safety domains
nor legal dynamics. The same construction applies to full monotone
interval branches elsewhere in the paper.

This example compares the two single-step tests for Lebesgue
probability. The next example addresses a different question: can a
product succeed when no choice of full-support reference probability
makes the single-step test work?

\subsection{Product contraction without a single-step reference test}
\label{subsec:product-example}
Single-step contraction can fail for every reference probability even
when a short product contracts. In the following example, densities
depend on one coordinate, so the product norm can be computed exactly.

Let $Q=(\mathbb Z/3\mathbb Z)^{\Nzero}$, let $P_u$ add $u$
to each coordinate, and use $g_u=P_u\sigma$ at period one. Put
\[
 K_0=[02]\cup[12]\cup[22],\quad K_1=[00]\cup[20],\quad
 K_2=[01]\cup[10]\cup[11]\cup[12]\cup[21].
\]
The only overlap is $[12]$. The cylinders $[02],[00],[01]$
require controls $0,1,2$, respectively, so $r=3$.
On $[12]$, choosing $0$ or $2$ changes the successor's first
letter from $2$ to $1$. Every Borel mixture on this entire overlap
is allowed. Illegal controls go to an isolated fixed cemetery point;
the resulting maps on $Q\sqcup\{\dagger\}$ are continuous.

For fair Bernoulli probability, the operators leave first-coordinate
functions invariant and have matrices $M_u/3$, where
\[
 M_0=\begin{pmatrix}0&0&0\\0&0&0\\1&1&1\end{pmatrix},\quad
 M_1=\begin{pmatrix}0&0&0\\1&0&1\\0&0&0\end{pmatrix},\quad
 M_2=\begin{pmatrix}1&1&1\\0&1&0\\0&1&0\end{pmatrix}.
\]
This follows from splitting the source measure by its first letter,
and retaining exactly the allowed two-letter cylinders. For a
positive product the full $L^\infty$ norm equals its norm on
$\one$, whose entire orbit stays in this invariant space. Thus this
finite representation computes the full product norms, not an
approximation to arbitrary Borel functions.

Let $A=M_0+M_1$, $B=M_0+M_2$, $C=M_1+M_2$.
The vectors after two of these unnormalized factors are
\[
\begin{array}{ccc}
(0,3,5)&(0,7,8)&(0,4,7)\\
(5,2,7)&(8,1,9)&(7,3,10)\\
(5,5,2)&(8,8,1)&(7,7,3).
\end{array}
\]
They are coordinatewise dominated by one of the vectors in
\[
 V=\{(0,7,8),(7,3,10),(7,7,3),(8,1,9),(8,8,1)\}.
\]
For $v=(p,q',r')$,
$Av=(0,p+r',p+q'+r')$,
$Bv=(p+q'+r',q',p+2q'+r')$, and
$Cv=(p+q'+r',p+q'+r',q')$.
For the five vectors of $V$, their largest resulting coordinate is,
respectively, $22,23,24,19,25$. Also $BCB\one=(17,8,25)$.
Consequently $b_3=25/27<1$. Submultiplicativity, rather than an
extrapolation from finite calculations, controls all lengths.
Theorem~\ref{thm:capacity-main} gives $M_1=\log3$, and every
unrefined feedback has that maximal-pattern rate. An explicit safe
selector chooses $0$ if $x_1=2$, otherwise $1$ if $x_1=0$ and
$x_0\ne1$, and otherwise $2$.

On the three target first-letter cylinders the densities are
$(0,0,1)$, $(0,2/3,1/3)$, and $(1,0,1/3)$.
Hence $\Theta_3=4/3$. This failure persists under \emph{every}
full-support probability $\nu$. If one of the safe branch measures
has no bounded $\nu$-density, the single-step criterion already
fails its prerequisites. Otherwise $g_0(K_0)\subset[2]$ and
$g_1(K_1)\subset[1]$ give disjoint density supports. At most two
densities are nonzero at a point, so their largest two sum to their
total. Its integral is
\[
 \sum_{u=0}^2\nu(K_u)=1+\nu([12])>1.
\]
Thus $\Theta_3(\nu)>1$ for every such $\nu$, including singular
or atomic probabilities. This excludes both original single-step
tests, but not all possible changes of presentation or iteration.

To compare global lists with state-dependent maximization, put
\[
 \Phi f=\max_{|S|=2}\sum_{u\in S}L_uf\qquad(f\ge0).
\]
On nonnegative functions the disjoint supports give
$\Phi=\sum_uL_u$. Under fair measure its first-letter matrix is
$\frac13\left(\begin{smallmatrix}1&1&1\\1&1&1\\1&2&1\end{smallmatrix}\right)$,
with positive eigenvector
$(1,1,(\sqrt{13}-1)/2)$ and eigenvalue
$(3+\sqrt{13})/6>1$. Hence $\|\Phi^n\one\|_\infty$ grows,
while every history of \emph{global} two-label lists decays.
Maximization by state and composition cannot be interchanged.


\section{Countable languages and common sampling}\label{sec:comb}

The capacity theorem chooses good observation positions separately
for each feedback. If only countably many name languages occur, they
can be served in turn by one increasing sequence. The main point is
to move on to later coordinates without translating a name language:
translation can lose patterns for a one-sided system. We instead
delete the old coordinates from a much larger maximizing window.

\subsection{Simultaneous attainment for pattern systems}

A \emph{finite-alphabet pattern system} is a nonempty set
$\Omega\subset E^{\Nzero}$ with $E$ finite and nonempty.  No topology or
shift invariance is imposed.  For finite $F\subset\Nzero$, write
\begin{align*}
 \Patt_\Omega(F)=\#\{x|_F:x\in\Omega\},
 \qquad
 p_\Omega^*(n)=
 \max_{\substack{F\subset\Nzero\\\#F=n}}\Patt_\Omega(F).
\end{align*}
With a fixed scale $\tau>0$, define
\begin{align*}
 \hpat(\Omega;\tau)&=\lim_{n\to\infty}
       \frac1{n\tau}\log p_\Omega^*(n),\\
 \hseq^A(\Omega;\tau)&=\limsup_{n\to\infty}
       \frac1{n\tau}\log
       \Patt_\Omega(\{a_1,\ldots,a_n\}).
\end{align*}
The maximum exists: its possible values form a nonempty subset of the finite
integer set $\{1,\ldots,(\#E)^n\}$.

We use $\Patt_\Omega(\varnothing)=1$ when needed.
Taking the product-topology closure of $\Omega$ changes none of its
finite projections: a cylinder meeting $\overline\Omega$ already meets
$\Omega$. Thus dispensing with closedness is only a matter of
formulation. The absence of shift invariance, in contrast, matters to the
placement of the coordinate sets.

\begin{lemma}[Pattern submultiplicativity and deletion]\label{lem:pattern-submult}
Let $\Omega\subset E^{\Nzero}$ be a pattern system and $q=\#E$.
For $m,n\in\N$ and finite $B\subset F\subset\Nzero$,
\begin{align}
 p_\Omega^*(m+n)&\le p_\Omega^*(m)p_\Omega^*(n),
 \label{eq:pattern-submult}\\
 \Patt_\Omega(F)&\le q^{\#(F\setminus B)}\Patt_\Omega(B).
 \label{eq:deletion} 
\end{align}
If $B\subset G$, then $\Patt_\Omega(B)\le\Patt_\Omega(G)$.
Consequently,
\begin{equation}\label{eq:fekete-pstar}
 \hpat(\Omega;\tau)
 =\inf_{n\ge1}\frac1{n\tau}\log p_\Omega^*(n),
 \qquad
 \log p_\Omega^*(n)\ge n\tau\hpat(\Omega;\tau).
\end{equation}
\end{lemma}

\begin{proof}
Split any $(m+n)$-element set into disjoint subsets of sizes $m$ and $n$.
Restriction sends a pattern on the union injectively to the pair of its two
restrictions, proving \eqref{eq:pattern-submult}.  Restriction from $F$ to
$B$ has fibers of cardinality at most $q^{\#(F\setminus B)}$, which proves
\eqref{eq:deletion}.  If $B\subset G$, every pattern on $B$ extends to a
pattern on $G$ using the same element of $\Omega$, giving monotonicity.
Fekete's lemma applied to $\log p_\Omega^*(n)$ proves
\eqref{eq:fekete-pstar}.
\end{proof}

For comparison, let $\Omega=\{0^\infty,10^\infty\}$. This set is a
closed one-sided subshift, but
$\Patt_\Omega(\{0\})=2$ and $\Patt_\Omega(\{m\})=1$ for $m\ge1$.
One cannot simply translate a maximizing set. Deleting a finite initial
interval has a different advantage: the number of lost coordinates is
bounded independently of the size of the maximizing set. Inequality
\eqref{eq:deletion} turns this bound into a controlled information loss.

\begin{theorem}[Simultaneous attainment]\label{thm:A}
Fix $\tau>0$.  For every countable family
$(\Omega_j)_{j\ge1}$ of finite-alphabet pattern systems, there is one
$A\in\mathscr S$ such that
\begin{equation}\label{eq:simultaneous-attainment}
 \hseq^A(\Omega_j;\tau)=\hpat(\Omega_j;\tau)
 \qquad(j\ge1).
\end{equation}
The sets $\Omega_j$ may be arbitrary nonempty subsets of their
finite-alphabet sequence spaces; shift invariance is not required.
\end{theorem}

\begin{proof}
At each stage we choose a large maximizing window for one language
and delete all positions preceding the end of the previous stage.
There are two losses to control: the deleted coordinates, and the
earlier observations still present in the denominator. Making the
new window large controls both. Each language receives infinitely
many stages, which suffices for its limsup.

Let $E_j$ be the alphabet of $\Omega_j$, put $q_j=\#E_j$ and
$h_j^*=\hpat(\Omega_j;\tau)$.  For every $A$ and $n$,
\begin{align*}
 \Patt_{\Omega_j}(\{a_1,\ldots,a_n\})\le p_{\Omega_j}^*(n),
\end{align*}
so \eqref{eq:fekete-pstar} gives
$\hseq^A(\Omega_j;\tau)\le h_j^*$.

Choose an enumeration of stages
\begin{align*}
 (j_k)_{k\ge1}=(1;\ 1,2;\ 1,2,3;\ 1,2,3,4;\ldots),
\end{align*}
where semicolons only separate finite groups. Every positive integer
occurs infinitely often. Choose $\eta_k>0$ with $\eta_k\to0$.  We recursively
construct nonempty finite sets $B_k$, each lying strictly to the right of all
previous ones.  Put
\begin{align*}
 G_{k-1}=\bigcup_{i<k}B_i,
 \qquad r_{k-1}=\#G_{k-1},
 \qquad M_{k-1}=\max G_{k-1},
\end{align*}
with $r_0=0$ and $M_0=-1$.

If $h_{j_k}^*>0$, choose $n_k$ so large that
\begin{equation}\label{eq:explicit-losses}
 n_k>M_{k-1}+1,
 \qquad
 \frac{(M_{k-1}+1)\log q_{j_k}}{n_k\tau}<\eta_k,
 \qquad
 \frac{r_{k-1}}{n_k}<\eta_k.
\end{equation}
Both $M_{k-1}$ and $r_{k-1}$ are fixed at this stage. The second
inequality pays for deleting up to $M_{k-1}+1$ positions; the third
makes the earlier $r_{k-1}$ observations negligible. Choose an
$n_k$-element $F_k$ attaining $p_{\Omega_{j_k}}^*(n_k)$ and set
\begin{align*}
 B_k=F_k\cap\{M_{k-1}+1,M_{k-1}+2,\ldots\}.
\end{align*}
At most $M_{k-1}+1$ coordinates are deleted, and the first condition in
\eqref{eq:explicit-losses} makes $B_k$ nonempty.  If $h_{j_k}^*=0$, set
$B_k=\{M_{k-1}+1\}$.

Let $A$ be the increasing enumeration of $\bigcup_{k\ge1}B_k$.  Fix $j$
with $h_j^*>0$ and restrict to the infinitely many stages with $j_k=j$.
Write $b_k=\#B_k$ and $R_k=r_{k-1}+b_k=\#G_k$; then $G_k$ consists of the
first $R_k$ entries of $A$.  By deletion, \eqref{eq:fekete-pstar}, and
monotonicity,
\begin{align}
 \log\Patt_{\Omega_j}(G_k)
 &\ge \log\Patt_{\Omega_j}(B_k)\notag\\
 &\ge \log p_{\Omega_j}^*(n_k)-(M_{k-1}+1)\log q_j\notag\\
 &\ge n_k\tau h_j^*-(M_{k-1}+1)\log q_j.
 \label{eq:endpoint-numerator}
\end{align}
Also $R_k\le r_{k-1}+n_k$.  At all sufficiently late such stages the last
line of \eqref{eq:endpoint-numerator} is positive.  Dividing by $R_k\tau$
and using \eqref{eq:explicit-losses} yields the explicit endpoint estimate
\begin{equation}\label{eq:endpoint-bound}
 \frac1{R_k\tau}\log\Patt_{\Omega_j}(G_k)
 >\frac{h_j^*-\eta_k}{1+\eta_k}.
\end{equation}
For a fixed $j$, the indices $k$ with $j_k=j$ tend to infinity, hence
$\eta_k\to0$ along those stages. Since every $B_i$ is nonempty,
$R_k\ge k\to\infty$. The limsup along these
endpoints is at least $h_j^*$, proving the lower
bound.  If $h_j^*=0$, the universal upper bound and nonnegativity give
equality for every $A$.  This proves \eqref{eq:simultaneous-attainment} for
all $j$.
\end{proof}

\begin{remark}[What is attained]\label{rem:attainment}
The theorem asserts equality of limsup rates. It does not assert a common
subsequence of endpoints on which every system is simultaneously close
to its maximal rate. The good endpoints used for different $j$ can be
different. Nor does it assert convergence of the normalized counts.
For example, take $\tau=1$, let $B\subset\Nzero$ and its complement both be infinite,
and let $\Omega_B$ consist of all binary sequences supported on $B$.
For any schedule, the two normalized logarithmic counts for $\Omega_B$
and $\Omega_{\Nzero\setminus B}$ add to $\log2$ at every $n$, although
both maximal-pattern rates are $\log2$. They cannot both converge to
$\log2$. The two limsups, however, can both equal $\log2$.

Exact finite-level maximization is also a stronger and generally false
requirement. Put
\begin{align*}
 \Omega=\{0x:x\in\{0,1\}^{\Nzero}\}\cup\{10^\infty,20^\infty\}
 \subset\{0,1,2\}^{\Nzero}.
\end{align*}
The unique maximizing singleton is $\{0\}$, with three patterns.
For a three-element set containing zero there are six patterns, whereas
any three positive coordinates realize eight patterns. No increasing
sequence maximizes the counts at both levels $n=1$ and $n=3$.
\end{remark}

\subsection{Partition names}\label{sec:partition-names}
To apply the preceding theorem to invariant partitions, use the atom
itineraries introduced in Section~\ref{sec:control}. We now allow an
arbitrary control set $U$; each partition still assigns only finitely
many blocks. Write a finite invariant partition as
$\C=(\{D_i:i\in I_\C\},\tau,\nu)$, with nonempty Borel atoms
$D_i$ and assigned blocks $\nu(D_i)\in U^\tau$ satisfying
\begin{equation}\label{eq:invariant-cover}
 \varphi(j,x,\nu(D_i))\in Q
       \qquad(x\in D_i,\ 0\le j\le\tau).
\end{equation}
Let $\mathfrak C_{\Bor}(Q,\tau)$ denote this class and
$\mathfrak C_{\clop}(Q,\tau)$ its clopen subclass. For variable
periods we use
\[
 \mathfrak C_{\Bor}(Q)=\bigcup_{\tau\ge1}\mathfrak C_{\Bor}(Q,\tau),
 \qquad \inf\varnothing=+\infty.
\]

The state map is $T_\C x=\varphi(\tau,x,\nu(D_i))$ on $D_i$.
Let $\iota_\C(x)$ be its atom itinerary and $W_N(\C;Q)$ the
set of length-$N$ prefixes. Thus $a_0\cdots a_{N-1}$ belongs to
$W_N(\C;Q)$ exactly when some $x\in Q$ satisfies
\begin{equation}\label{eq:admissible-word}
 \varphi(j\tau,x,\nu(D_{a_0})\cdots\nu(D_{a_{N-1}}))
       \in D_{a_j}\qquad(0\le j<N).
\end{equation}
These are Kawan's admissible words \cite[Definition~2.9]{Kawan2013};
the induced map is also used in \cite[Section~2]{TomarKawanZamani2022}.
Condition~\eqref{eq:invariant-cover} ensures safety between successive
block boundaries. After the last block the same rule supplies the next
one, so every finite witness continues indefinitely. A suffix is
witnessed by the state at its starting block. In particular,
\begin{equation}\label{eq:word-submult}
 |W_{M+N}(\C;Q)|\le |W_M(\C;Q)|\,|W_N(\C;Q)|.
\end{equation}

We shall sometimes use the compact name space
\begin{equation}\label{eq:name-system}
 \Sigma_\C=\{a\in I_\C^{\Nzero}:
       a|_{[0,N-1]}\in W_N(\C;Q)\text{ for every }N\}.
\end{equation}
It is closed, and suffix closure gives
$\sigma\Sigma_\C\subset\Sigma_\C$. The cylinder description yields
\begin{equation}\label{eq:itinerary-closure}
 \Sigma_\C=\overline{\iota_\C(Q)}.
\end{equation}
Every cylinder meeting this closure already contains an itinerary.
Hence its finite patterns are exactly the actual patterns, and our
earlier notation becomes
\begin{align}
 \Patt_\C(F)&=|\{a|_F:a\in\Sigma_\C\}|,
       \label{eq:control-pattern-count}\\
 p_\C^*(n)&=\max_{|F|=n}\Patt_\C(F),\label{eq:pstar-control}\\
 h^*(\C;Q)&=h^*(\Sigma_\C;\tau),\label{eq:hpat-control}\\
 h_{\rm seq}^A(\C;Q)&=h_{\rm seq}^A(\Sigma_\C;\tau).
       \label{eq:hseq-control}
\end{align}
We also need the consecutive-name rate
\begin{equation}\label{eq:ordinary-partition-entropy}
 h(\C;Q)=\lim_N\frac{\log|W_N(\C;Q)|}{N\tau},
\end{equation}
whose limit follows from \eqref{eq:word-submult}.

\begin{remark}[Infinite names in the closure]\label{rem:witness}
For clopen atoms the maps $T_\C$ and $\iota_\C$ are continuous,
so compactness gives $\Sigma_\C=\iota_\C(Q)$. With Borel atoms,
some infinite names in the closure can lack a single state witness.
For example, take $Q=[0,1]$, one control $f(x)=\min\{2x,1\}$,
and atoms $D_0=(0,1/2)$, $D_1=Q\setminus D_0$. The state
$2^{-N-1}$ realizes $0^N$, but $0^\infty$ has no witness:
zero is not in $D_0$, and each positive point eventually leaves it.
The finite counts in \eqref{eq:control-pattern-count} are nevertheless
actual trajectory counts, by the cylinder argument above.
\end{remark}

\begin{corollary}[Fixed-language attainment]\label{cor:fixed}
For every finite-alphabet pattern system $\Omega$ and every $\tau>0$,
\begin{equation}\label{eq:fixed-attainment}
 \hpat(\Omega;\tau)=\max_{A\in\mathscr S}\hseq^A(\Omega;\tau).
\end{equation}
The same identity holds for every finite Borel invariant partition $\C$ after
putting $\Omega=\Sigma_\C$.
\end{corollary}


\subsection{Feedback classes with countably many languages}\label{sec:minimax}

Fix a compact controlled-invariant set $Q$ and one period $\tau\in\N$.
For two invariant partitions $\C$ and $\D$ of this period, write
$\C\sim_{\lang}\D$ if there is a bijection
$\pi:I_\C\to I_\D$ such that, for every $N\in\N$,
\begin{equation}\label{eq:language-equivalence}
 \pi^N(W_N(\C;Q))=W_N(\D;Q),
\end{equation}
where $\pi^N$ acts coordinatewise.  This is \emph{language equivalence}.
Identity, inverse, and composition of the alphabet bijections show that
$\sim_{\lang}$ is an equivalence relation.  If $F\subset\Nzero$ is nonempty and finite and
$m=\max F$, then \eqref{eq:language-equivalence} at length $m+1$, followed by
projection onto $F$, gives a bijection between the two $F$-pattern sets.
Thus language-equivalent partitions have identical finite-coordinate pattern
counts and identical sequence and maximal-pattern complexities.

Fix a nonempty class $\mathfrak C\subset\mathfrak C_{\Bor}(Q,\tau)$.
Every individual rate is nonnegative and finite, so its infimum over
this class is finite even without a uniform bound on the numbers of atoms.

The simultaneous-attainment theorem applies once there are only
countably many such languages, up to renaming their letters. Taking
infima along its common sequence gives the reverse of the elementary
minimax inequality. This yields the following direct corollary.

\begin{corollary}[Minimax over countably many language types]\label{thm:B}
Let $\varnothing\ne\mathfrak C\subset\mathfrak C_{\Bor}(Q,\tau)$ and suppose
that the quotient $\mathfrak C/{\sim_{\lang}}$ is finite or countable.  There
is one $A\in\mathscr S$ such that
\begin{equation}\label{eq:pointwise-class-attainment}
 \hseq^A(\C;Q)=\hpat(\C;Q)
 \qquad(\C\in\mathfrak C).
\end{equation}
Consequently,
\begin{equation}\label{eq:minimax}
 \max_{B\in\mathscr S}\inf_{\C\in\mathfrak C}
 \hseq^B(\C;Q)
 =\inf_{\C\in\mathfrak C}\hpat(\C;Q).
\end{equation}
The outer maximum is attained. The pattern infimum and the inner infimum at the common attaining schedule are minima, by Corollary~\ref{cor:pattern-minimum}.
\end{corollary}

\begin{proof}
Choose one representative $\C_j$ of each language-equivalence class.
Enumerate the representatives, repeating them periodically when the quotient
is finite, and apply Theorem~\ref{thm:A} to the resulting countable family
$\Sigma_{\C_j}$, using the common scale $\tau$.  The resulting $A$ gives equality for each
representative.  Language equivalence preserves all finite pattern counts, so
it gives \eqref{eq:pointwise-class-attainment} for every $\C\in\mathfrak C$.
Taking the infimum over $\C$ gives equality at this $A$.

For any $B\in\mathscr S$, Corollary~\ref{cor:fixed} gives
$\hseq^B(\C;Q)\le\hpat(\C;Q)$ for every $\C$.  Taking first the infimum over
$\C$ and then the supremum over $B$ proves the reverse minimax inequality and
hence \eqref{eq:minimax}.  If the pattern infimum is zero, the pointwise upper
bound and nonnegativity in fact give
$\inf_{\C\in\mathfrak C}\hseq^B(\C;Q)=0$ for every $B\in\mathscr S$, without
any countability assumption.
\end{proof}

Language equivalence allows alphabet renumbering while preserving all
finite pattern counts. It does not merge distinct atoms merely because
their assigned controls agree. The countability assumption concerns
the languages of the chosen partition class at the fixed period.

\begin{proposition}[Countability of clopen feedbacks]
\label{prop:digital}
Assume that $Q$ is compact metrizable and that the control alphabet $U$ is
countable.  Then
$\mathfrak C_{\clop}(Q,\tau)$ is countable.  Hence, if this class is nonempty,
Corollary~\ref{thm:B} applies to it.  In particular this holds when $U$ is a
finite digital alphabet.
\end{proposition}

\begin{proof}
A compact metrizable space has a countable base.  Every clopen subset is
compact and hence is a finite union of basic open sets contained in it.  Thus
$Q$ has only countably many clopen subsets and only countably many finite
clopen partitions.  Since $\tau$ is finite, $U^\tau$ is countable.  A finite
partition admits only countably many block assignments.  The invariant assignments form a
subcollection.
\end{proof}

Corollary~\ref{thm:B} fixes $\tau$.  It does not produce one common
physical-time schedule for partitions of different periods and proves no
minimax identity over $\mathfrak C_{\Bor}(Q)$.



The countability hypothesis has concrete limitations.
Appendix~\ref{app:pointwise} develops the elementary obstruction
of Example~\ref{ex:uncountable} into autonomous constructions.
In particular, Proposition~\ref{prop:finite-pointwise}
uses four distinct controls and has $L_1=M_1=\log2$, although no
sequence attains every feedback's own maximal-pattern rate. It therefore
separates pointwise common attainment from a common full-class minimum.


\section{Comparison with classical and invariance entropy}\label{sec:comparison}

A fixed-period pattern rate can be positive when invariance entropy
vanishes. We first place the counts beside the usual entropy notions
and explain what happens when the control period varies. A symbolic
realization then gives an explicit separation at period one.

\subsection{Classical entropy and sequence entropy}\label{sec:background}
For a continuous map $T:Z\to Z$ of a compact metric space and a
finite open cover $\mathcal V$, replace the name count by
\[
 c_{\mathcal V}(F)=N\left(\bigvee_{j\in F}T^{-j}\mathcal V\right),
\]
where $N$ denotes minimum subcover cardinality. Consecutive windows,
prefixes of $A$, and maximizing $n$-element windows give, respectively,
\begin{align*}
 h_{\rm top}(T,\mathcal V)
   &=\lim_n n^{-1}\log c_{\mathcal V}(\{0,\ldots,n-1\}),\\
 h_{\rm top}^A(T,\mathcal V)
   &=\limsup_n n^{-1}\log c_{\mathcal V}(\{a_1,\ldots,a_n\}),\\
 h_{\rm top}^*(T,\mathcal V)
   &=\lim_n n^{-1}\log\max_{|F|=n}c_{\mathcal V}(F).
\end{align*}
Submultiplicativity gives the first and third limits. Taking the
supremum of the first over finite open covers defines
$h_{\rm top}(T)$.

For homeomorphisms, Huang and Ye prove
\begin{equation}\label{eq:HY-background}
 h_{\rm top}^*(T,\mathcal V)
   =\max_A h_{\rm top}^A(T,\mathcal V)
\end{equation}
\cite[Theorem~2.1(1)]{HuangYe2009}. Their proof places successive
maximizing windows farther to the right by translation. Surjectivity
preserves cover counts under that translation. The deletion argument
in Section~\ref{sec:comb} instead applies to arbitrary name sets,
including those of non-surjective maps.

For a clopen partition, $c_{\mathcal V}(F)$ is the number of
nonempty atoms of the join, so it is precisely a name count. For an
overlapping cover it is a minimum covering number, as illustrated
by Corollary~\ref{cc-cor:covergap}. A further distinction is needed
for measurable entropy: Shannon entropy weights the atoms by an
invariant probability, whereas our counts retain every nonempty
atom. Example~\ref{ex:periodic-support} shows that even a countable
set of initial states can change the counting rate.

On a subshift the coordinate partition gives the name counts used
here. For ordinary entropy its consecutive joins refine every fixed
cylinder partition without changing the exponential rate. Sparse
observations behave differently: observing a window of several
coordinates at each selected time can increase the rate per
observation. Accordingly, we write ``$\mathrm{coord}$'' for the
coordinate rates below; we do not take a supremum over open covers.

\subsection{Comparison with Kawan's invariance entropy}
For every invariant partition, consecutive coordinates are one permissible
choice in the maximal-pattern count. Therefore
\begin{equation}\label{eq:ordinary-pattern-one-sided}
 h(\C;Q)\le\hpat(\C;Q)\le\frac{\log\#I_\C}{\tau}.
\end{equation}
Kawan's Theorem~2.3, with natural logarithms, gives
\begin{equation}\label{eq:kawan-characterization}
 \hinv(Q)=\inf_{\C\in\mathfrak C_{\Bor}(Q)}h(\C;Q).
\end{equation}
Its hypotheses are satisfied by the full-control discrete-time systems
used here: $X$ is metrizable, each controlled map is continuous, and
$Q$ is nonempty, compact and controlled invariant.
The period in this formula is allowed to vary.

The comparison uses two constructions: a spanning family gives a
finite invariant partition, and the words of an invariant partition
give spanning families at longer horizons. First, spanning families
concatenate, so
$r_{\rm inv}(m+n,Q)\le r_{\rm inv}(m,Q)r_{\rm inv}(n,Q)$ when the
right-hand side is finite. Moreover $r_{\rm inv}(n,Q)$ is nondecreasing.
If any finite-horizon spanning number is finite, then $r_{\rm inv}(1,Q)$
is finite, all of them are finite by concatenation, and Fekete's lemma
gives a limit in \eqref{eq:hinv-definition}. If none is finite, that
limit equals $+\infty$.

Let $S_N=\{\omega^1,\ldots,\omega^m\}$ be a minimal finite
$N$-step spanning family. Its safe sets are
\begin{align*}
 K_i=Q\cap\bigcap_{s=1}^{N}
       \{x\in X:\varphi(s,x,\omega^i)\in Q\}.
\end{align*}
They are closed in $Q$, since finite control compositions are continuous
and $Q$ is closed. They cover $Q$. Set
$D_1=K_1$ and $D_i=K_i\setminus\bigcup_{j<i}K_j$, discard empty sets,
and assign $\omega^i$ to $D_i$. This gives a Borel invariant partition $\C_N$ of period $N$
with at most $m$ atoms, using the spanning construction in Kawan's
Theorem~2.3.

For the other inequality in \eqref{eq:kawan-characterization}, all
admissible $k$-block words of any $\C$ form a spanning family after their
assigned controls are concatenated. Indeed, from each initial state the
feedback rule selects one such admissible word. Consequently
$r_{\rm inv}(k\tau,Q)\le\#W_k(\C;Q)$, and taking limits gives
$\hinv(Q)\le h(\C;Q)$. The reverse inequality follows from
$h(\C_N;Q)\le N^{-1}\log r_{\rm inv}(N,Q)$ and $N\to\infty$.
If no finite spanning family exists, no finite invariant partition exists,
and both sides are $+\infty$.

Thus, for every nonempty fixed-period subclass
$\mathfrak C\subset\mathfrak C_{\Bor}(Q,\tau)$, the general comparison is
\begin{equation}\label{eq:three-quantities}
 \hinv(Q)
 \le\inf_{\C\in\mathfrak C}h(\C;Q)
 \le\inf_{\C\in\mathfrak C}\hpat(\C;Q).
\end{equation}
The Thue--Morse realization below shows that the last quantity can
strictly exceed $\hinv(Q)$.

\subsection{Why the unrestricted all-period infimum collapses}
Allowing the control period to vary gives the quantity
\begin{align*}
 H_{\rm pat,all}^{*,\Bor}(Q)
 :=\inf_{\tau\in\N}\inf_{\C\in\mathfrak C_{\Bor}(Q,\tau)}
       \hpat(\C;Q).
\end{align*}
Kawan's spanning construction gives the following direct consequence.
The alphabet of an $N$-step partition already bounds every sampled
count, regardless of the observation gaps.

\begin{proposition}[All-period collapse]\label{prop:all-period-collapse}
In the full-control discrete-time setting of Section~\ref{sec:control}, every
nonempty compact controlled-invariant set $Q$ satisfies
\begin{equation}\label{eq:all-period-collapse}
 H_{\rm pat,all}^{*,\Bor}(Q)=\hinv(Q).
\end{equation}
\end{proposition}

\begin{proof}
The lower bound follows from
\eqref{eq:three-quantities}. For the upper bound use the same
$\C_N$ just constructed. Since its alphabet has at most
$r_{\rm inv}(N,Q)$ symbols, every pattern on $k$ coordinates has at most
$r_{\rm inv}(N,Q)^k$ possibilities. Hence
\begin{align*}
 H_{\rm pat,all}^{*,\Bor}(Q)
 \le\hpat(\C_N;Q)
 \le\frac{\log r_{\rm inv}(N,Q)}{N}.
\end{align*}
Let $N\to\infty$. The infinite-entropy case was already covered by
nonexistence of finite invariant partitions. This proves
\eqref{eq:all-period-collapse}.
\end{proof}

The all-period identity does not produce a common physical observation
sequence. A block index $a_i$ corresponds to time $\tau a_i$, which
changes with $\tau$. It says only that packing longer control words
into one observed symbol can lower the normalized pattern rate.

\subsection{Symbolic safety realizations}\label{sec:realization}

To transfer a difference between consecutive and sparse symbolic
counts to a control problem, we force the safe control to equal the
current state symbol. This leaves no choice of physical label, while
still allowing arbitrary Borel refinements of its observation.

\subsubsection{Matching the control to the current symbol}

Let $E$ be a finite nonempty discrete alphabet and let
$\varnothing\ne Y\subset E^{\Nzero}$ be closed with $\sigma(Y)\subset Y$.
Write
\begin{align*}
 \Lcal_N(Y)=\{y_0\cdots y_{N-1}:y\in Y\},
 \qquad
 P_Y(F)=\#\{y|_F:y\in Y\}.
\end{align*}
The coordinate complexities are
\begin{align*}
 h_{\rm seq,coord}^A(Y)
 &=\limsup_{n\to\infty}\frac1n
   \log P_Y(\{a_1,\ldots,a_n\}),\\
 h_{\rm pat,coord}^*(Y)
 &=\lim_{n\to\infty}\frac1n
   \log\max_{\substack{F\subset\Nzero\\\#F=n}}P_Y(F).
\end{align*}

Adjoin an isolated cemetery point and give
$X=Y\sqcup\{\dagger\}$ the topological disjoint-union topology.  Put $U=E$
and define
\begin{equation}\label{eq:matching-control}
 f(y,u)=
 \begin{cases}
   \sigma y,&y_0=u,\\
   \dagger,&y_0\ne u,
 \end{cases}
 \qquad
 f(\dagger,u)=\dagger.
\end{equation}

\begin{proposition}[Symbolic-to-control realization]\label{thm:C}
The matching-control system \eqref{eq:matching-control} is continuous
on compact metrizable $X$, and $Y$ is controlled invariant. Its
finite-horizon spanning count is, for every $N\in\N$,
\begin{equation}\label{eq:rinv-language}
 r_{\rm inv}(N,Y)=\#\Lcal_N(Y),
\end{equation}
and therefore
\begin{equation}\label{eq:hinv-htop}
 \hinv(Y)=\htop(Y,\sigma)
 =\lim_{N\to\infty}\frac1N\log\#\Lcal_N(Y).
\end{equation}
At period one identify $U^1$ with $U$. If
$\C=(\{D_i:i\in I_\C\},1,\nu)\in\mathfrak C_{\Bor}(Y,1)$, set
$\lambda_\C(i)=\nu(D_i)\in E$.  Then every atom satisfies
\begin{equation}\label{eq:atom-cylinder}
 D_i\subset[\lambda_\C(i)]_Y
 :=\{y\in Y:y_0=\lambda_\C(i)\}.
\end{equation}
Moreover, the one-block label map
\begin{equation}\label{eq:label-factor}
 \Lambda_\C:\Sigma_\C\longrightarrow Y,
 \qquad (a_j)_{j\ge0}\longmapsto(\lambda_\C(a_j))_{j\ge0},
\end{equation}
is a continuous, onto, shift-commuting map.
Consequently, for every $A\in\mathscr S$,
\begin{align}
 \inf_{\C\in\mathfrak C_{\Bor}(Y,1)}\hseq^A(\C;Y)
 &=\inf_{\C\in\mathfrak C_{\clop}(Y,1)}\hseq^A(\C;Y)
 =h_{\rm seq,coord}^A(Y),
 \label{eq:seq-realization}\\
 \inf_{\C\in\mathfrak C_{\Bor}(Y,1)}\hpat(\C;Y)
 &=\inf_{\C\in\mathfrak C_{\clop}(Y,1)}\hpat(\C;Y)
 =h_{\rm pat,coord}^*(Y),
 \label{eq:pat-realization}\\
 \inf_{\C\in\mathfrak C_{\Bor}(Y,1)}h(\C;Y)
 &=\inf_{\C\in\mathfrak C_{\clop}(Y,1)}h(\C;Y)
 =\htop(Y,\sigma).
 \label{eq:ordinary-realization}
\end{align}
Every Borel and clopen infimum in
\eqref{eq:seq-realization}--\eqref{eq:ordinary-realization} is attained by
the same first-coordinate cylinder partition.
\end{proposition}

\begin{proof}
Safety forces the control word to match the state prefix. This gives
the spanning count directly. For a refined partition, we then project
atom labels back to those forced controls, obtaining the lower bound
for every finite observation window. The first-coordinate partition
attains all these bounds.

The space $X$ is the union of the two clopen compact sets $Y$ and
$\{\dagger\}$, hence is compact metrizable.  For fixed $u\in E$, the sets
$[u]_Y$, $Y\setminus[u]_Y$, and $\{\dagger\}$ are clopen; on them $f_u$ is,
respectively, the shift or a constant.  Thus $f_u$ is continuous.  Since $U$
is finite discrete, $f$ is jointly continuous.  Choosing $u=y_0$ sends $y$
to $\sigma y\in Y$, proving controlled invariance.

A direct induction gives
\begin{equation}\label{eq:safe-iff-match}
 \varphi(n,y,\omega)\in Y\ (0\le n\le N)
 \quad\Longleftrightarrow\quad
 \omega_j=y_j\ (0\le j<N).
\end{equation}
Indeed, after $j$ matches the state is $\sigma^j y$, whereas a mismatch enters
$\dagger$ permanently.  Hence a length-$N$ control word is safe from at
least one state of $Y$ exactly when it belongs to $\Lcal_N(Y)$, and then it
safely covers precisely its nonempty prefix cylinder.  Distinct language
words cover disjoint cylinders, and every such cylinder is nonempty.
A spanning family must therefore contain every word of $\Lcal_N(Y)$;
those words also suffice. This proves \eqref{eq:rinv-language}.  The language counts are
submultiplicative, so their normalized logarithms converge.

For completeness, this limit is the usual topological entropy even when
$\sigma|_Y$ is not onto.  The join of the first-coordinate cylinder cover
through times $0,\ldots,N-1$ has exactly $\#\Lcal_N(Y)$ nonempty atoms.
Conversely, given a finite open cover of compact $Y$, the partition into
cylinders of some fixed length $m$ refines that cover; its $N$-fold dynamical join is
then refined by the length-$(N+m-1)$ cylinders.  Taking exponential rates proves
\eqref{eq:hinv-htop}.

Let $\C=(\{D_i:i\in I_\C\},1,\nu)$ be a Borel invariant partition.  If
$D_i$ contained $y$ with $y_0\ne\lambda_\C(i)$, its assigned control would send $y$ to
$\dagger\notin Y$, contradicting safety.  This proves
\eqref{eq:atom-cylinder}.

Take $a=(a_j)_{j\ge0}\in\Sigma_\C$ and put $z_j=\lambda_\C(a_j)$.  For each
$N$, choose a witness $y^{(N)}$ for the first $N$ symbols of $a$.  Induction
using matching controls and \eqref{eq:atom-cylinder} gives
$y^{(N)}_j=\lambda_\C(a_j)=z_j$ for $0\le j<N$.  Hence
$y^{(N)}\to z$ in the product topology; since $Y$ is closed, $z\in Y$.  Thus
$\Lambda_\C$ is well defined, and it is continuous: the preimage of a cylinder in $E^{\Nzero}$ is a finite
union of cylinders in $I_\C^{\Nzero}$, restricted to $\Sigma_\C$.
It is shift commuting because
$(\Lambda_\C(\sigma a))_j=\lambda_\C(a_{j+1})
=(\sigma\Lambda_\C(a))_j$.  Given $y\in Y$, let $a_j$ be the unique index with
$\sigma^j y\in D_{a_j}$.  Its assigned label is $y_j$, and $y$ witnesses
every finite prefix of this atom itinerary.  Thus \eqref{eq:label-factor} is onto.  The factor maps compact name spaces: the limiting point $z$ need
not follow the same Borel atom itinerary $a$. Only its forced control
labels are used in the convergence argument.

For every finite $F$, the coordinate projection of $\Lambda_\C$ is onto, so
\begin{equation}\label{eq:partition-dominates-coordinate}
 \Patt_\C(F)\ge P_Y(F).
\end{equation}
It follows that every Borel, and hence every clopen, invariant partition has
sequence, maximal-pattern, and consecutive-name entropy at least the
corresponding coordinate quantity.

Let $E_Y=\{e\in E:[e]_Y\ne\varnothing\}$ and define
\begin{align*}
\C_0=(\{D_e=[e]_Y:e\in E_Y\},1,\nu_0),
 \qquad \nu_0(D_e)=e.
\end{align*}
This is a clopen invariant partition.  Under the symbol identification,
$W_N(\C_0;Y)=\Lcal_N(Y)$ and $\Sigma_{\C_0}=Y$; the latter follows from
closedness of $Y$.  Equality holds in \eqref{eq:partition-dominates-coordinate}
for every $F$.  Thus $\C_0$ attains all the lower bounds, proving
\eqref{eq:seq-realization}--\eqref{eq:ordinary-realization}.
\end{proof}


\subsection{A Thue--Morse separation}\label{sec:thue-morse}

Let $\vartheta(0)=01$, $\vartheta(1)=10$, and let
\begin{align*}
 \mathbf t=\lim_{r\to\infty}\vartheta^r(0)
 =0110100110010110\cdots.
\end{align*}
Equivalently, $t_n=s_2(n)\pmod2$, where $s_2(n)$ is the number of ones in
the binary expansion of $n$.  Let
\begin{equation}\label{eq:thue-morse-shift}
 Y=\overline{\{\sigma^m\mathbf t:m\ge0\}}
 \subset\{0,1\}^{\Nzero}.
\end{equation}
This is the one-sided Thue--Morse orbit closure.  The following direct computation supplies both rates used below.

\begin{proposition}[Exact symbolic rates]\label{prop:thue-morse}
For every $N\ge1$,
\begin{equation}\label{eq:TM-linear}
 \#\Lcal_N(Y)\le8N.
\end{equation}
Hence $\htop(Y)=0$.  On the single schedule
\begin{equation}\label{eq:TM-schedule}
 A_{\rm TM}=(a_i)_{i\ge1},
 \qquad a_i=4^{i-1}-1,
\end{equation}
one has, for every $k\ge1$,
\begin{equation}\label{eq:TM-full-pattern}
 P_Y(\{a_1,\ldots,a_k\})=2^k.
\end{equation}
The chosen coordinate set contains zero.  Since a cylinder is clopen, its
patterns in the orbit closure are exactly the patterns obtained along shifts
of $\mathbf t$; this calculation uses actual shifts of the fixed point.  Consequently
\begin{equation}\label{eq:TM-complexities}
 h_{\rm seq,coord}^{A_{\rm TM}}(Y)
 =h_{\rm pat,coord}^*(Y)=\log2.
\end{equation}
\end{proposition}

\begin{proof}
A finite word is in $\Lcal_N(Y)$ exactly when it occurs as a factor of
$\mathbf t$: one direction follows from the orbit, and the other from
convergence in a prefix cylinder.  Choose $\ell=\lceil\log_2N\rceil$.  The
identity $\mathbf t=\vartheta^\ell(\mathbf t)$ decomposes $\mathbf t$ into
aligned supertiles of length $2^\ell$.  Since $N\le2^\ell$, every length-$N$
factor is contained in $\vartheta^\ell(ab)$ for an adjacent two-letter factor
$ab$.  There are at most four choices for $ab$ and at most $2^\ell$ starting
offsets in the first supertile.  Therefore
$\#\Lcal_N(Y)\le4\,2^\ell\le8N$, proving zero topological entropy.

For the sparse lower bound, we encode each prescribed bit at an even
binary position. The intervening odd positions stop carries, while
one higher bit adjusts the total parity.

Fix $k\ge1$ and
$\varepsilon=(\varepsilon_0,\ldots,\varepsilon_{k-1})\in\{0,1\}^k$.
Put
\begin{align*}
 r=\#\{0\le i<k:\varepsilon_i=0\},
 \qquad
 q=\sum_{\substack{0\le i<k\\\varepsilon_i=0}}4^i
 +(r\bmod2)4^k.
\end{align*}
Only even binary positions occur in this sum.  The binary digit sum of $q$ is
$r+(r\bmod2)$, which is even, so $t_q=0$.  If $\varepsilon_i=1$, bit $2i$
of $q$ is zero, and adding $4^i=2^{2i}$ changes the digit-sum parity.  If
$\varepsilon_i=0$, bit $2i$ is one and bit $2i+1$ is zero; adding $2^{2i}$
replaces these two bits by zero and one without a further carry, so the parity
is unchanged.  In both cases
\begin{equation}\label{eq:TM-digit-identity}
 t_{q+4^i}=\varepsilon_i\qquad(0\le i<k).
\end{equation}
The extra bit at position $2k$ only corrects the initial parity of $q$.
For $i<k$, a carry stops at the empty odd position $2i+1$, which is
strictly below $2k$. The parity correction therefore interferes with
none of these tests.

Set $n=q+1$. Since $a_{i+1}=4^i-1$,
\begin{align*}
 (\sigma^n\mathbf t)_{a_{i+1}}
 =t_{q+4^i}=\varepsilon_i.
\end{align*}
Thus every binary $k$-pattern occurs on the first $k$ members of
$A_{\rm TM}$, proving \eqref{eq:TM-full-pattern}.  The binary alphabet gives
the reverse upper bound $p_Y^*(k)\le2^k$, and
\eqref{eq:TM-complexities} follows.
\end{proof}

\begin{corollary}[Zero invariance entropy, full binary sparse-time complexity]
\label{cor:separation}
Apply the matching-control construction \eqref{eq:matching-control} to the
Thue--Morse subshift \eqref{eq:thue-morse-shift}.  Then
\begin{equation}\label{eq:strict-control-separation}
 \hinv(Y)=0
 <\log2
 =\inf_{\C\in\mathfrak C_{\Bor}(Y,1)}\hpat(\C;Y)
 =\inf_{\C\in\mathfrak C_{\clop}(Y,1)}\hpat(\C;Y).
\end{equation}
Moreover, the explicit schedule $A_{\rm TM}$ satisfies
\begin{align}
 \max_{A\in\mathscr S}\inf_{\C\in\mathfrak C_{\Bor}(Y,1)}
 \hseq^A(\C;Y)
 &=\inf_{\C\in\mathfrak C_{\Bor}(Y,1)}\hseq^{A_{\rm TM}}(\C;Y)
 =\log2,
 \label{eq:TM-Borel-sequence}\\
 \max_{A\in\mathscr S}\inf_{\C\in\mathfrak C_{\clop}(Y,1)}
 \hseq^A(\C;Y)
 &=\inf_{\C\in\mathfrak C_{\clop}(Y,1)}\hseq^{A_{\rm TM}}(\C;Y)
 =\log2.
 \label{eq:TM-clopen-sequence}
\end{align}
\end{corollary}

\begin{proof}
Combine Proposition~\ref{thm:C} with Proposition~\ref{prop:thue-morse}.  For any
schedule the binary coordinate complexity is at most $\log2$, while
$A_{\rm TM}$ attains $\log2$.  Equations
\eqref{eq:seq-realization}--\eqref{eq:pat-realization} give the Borel and
clopen infima.
\end{proof}

The separation persists after minimization over all Borel invariant
partitions at period one. It therefore distinguishes the fixed-period
pattern infimum from invariance entropy.


For the matching-control Thue--Morse example, the collapse is particularly
visible. At period $N$, the length-$N$ cylinder partition with its matching
control blocks is clopen, has at most $8N$ atoms and has pattern rate at
most $N^{-1}\log(8N)$. Thus even the all-period clopen infimum is zero
in this example, while the period-one Borel infimum is $\log2$.
Changing the control period changes which information is packed into
one observed symbol and the cost assigned to that symbol.

\section{Affine examples and feedback consistency}\label{sec:affine}

For expanding interval branches, a prescribed word can be followed
only from a small set of initial states. Covering the interval therefore
requires many words. The first example matches this volume bound by
a feedback that eventually enters an invariant core. The second uses
protected inverse branches to force every sampled word under every
feedback. The four-branch example then separates what these arguments
can establish: its maximal-pattern minimum is $\log2$, while its
common-sampling value is still only bounded. We finish by showing why
an efficient open-loop covering need not define a feedback.

\subsection{Integer affine branches and an absorbing feedback}
Let $m\ge2$, $D\ge m$ be integers, $R=D/(m-1)$,
$Q=[0,R]$, and controls $u=0,\ldots,D$ act legally by $mx-u$.
On $X=[-1,R+1]$ define the continuous ambient maps
\[
 f_u(x)=\max\{-1,\min\{R+1,mx-u\}\}.
\]
For $x\in Q$, this is in $Q$ exactly when $mx-u$ is in $Q$,
and then truncation has no effect. The safety intervals
$K_u=[u/m,(u+R)/m]$ cover $Q$. Fix any period $\tau\ge1$.

\begin{proposition}\label{prop:integer}
For this system, over all legal Borel block feedbacks and finite
Borel atom refinements,
\[
 \min_s h_{(0,1,\ldots)}(s)=L_\tau=M_\tau=\log m.
\]
There is an actual Borel selector $s_\tau$ with
$h_A(s_\tau)=h^*(s_\tau)=\log m$ for every infinite increasing
$A$, although its absorption time has no uniform bound.
\end{proposition}
\begin{proof}
Expansion gives the lower bound by making each word's initial-state
fiber small. For the upper bound we construct a rule under which
every state except the right endpoint eventually enters the usual
$m$-digit core. Entry times are not uniformly bounded,
so we count where the transition can occur among the observed times,
rather than bounding the elapsed waiting time.

For any actual length-$n$ block word, concatenate its controls into
$w_0\cdots w_{n\tau-1}$. Safety gives
\[
 x_{n\tau}=m^{n\tau}x-\sum_{j=0}^{n\tau-1}
                   w_jm^{n\tau-1-j}\in[0,R].
\]
Its initial-state fiber is contained in an interval of length
$R/m^{n\tau}$, whether or not that fiber itself is an interval.
These containing intervals cover $Q$, so there are at least
$m^{n\tau}$ actual consecutive block words. This gives the
full-class lower bound at the common sequence $(0,1,\ldots)$.

Define the one-step rule $v(x)=\min\{D,\lfloor mx\rfloor\}$
and $T(x)=mx-v(x)$. If $x<D/m$, the next state is in $[0,1)$;
if $x\ge D/m$, it is $mx-D\in[0,R]$.
The core $[0,1)$ is invariant and uses only
$\Gamma=\{0,\ldots,m-1\}$. Both $0$ and $R$ are fixed;
$D/m$ uses $D$ and goes to zero.
For $x<R$, while $D$ is used the state is
$R-m^j(R-x)$. Since $R-x>0$, eventually it is below $D/m$,
then one transitional control $v<D$ sends it into the core.
Only $R$ uses $D$ forever. The absorption times are unbounded
as $x$ tends to $R$.

From the current state define
$s_\tau(x)=(v(x),v(Tx),\ldots,v(T^{\tau-1}x))$.
It is Borel, each intermediate step is safe, and its actual endpoint
map is $T^\tau$. Put $\mathcal B=\Gamma^\tau$,
$b=m^\tau$, and $d_0=(D,\ldots,D)$.
Its block names are contained in
\[
 \bigcup_{j\ge0}d_0^j\mathcal C_\tau\mathcal B^{\Nzero}
 \ \cup\ \{d_0^\infty\},\qquad
 |\mathcal C_\tau|\le D\frac{m^\tau-1}{m-1}=:C_\tau.
\]
The transitional block has some initial $D$'s, one control below
$D$, and then core controls; this proves the stated count.
For any window of $n$ block positions, an unobserved transition
leaves just one of $n+1$ cuts between constant initial labels and
core labels. A sampled transition has $n$ possible locations and
at most $C_\tau$ labels. Thus
\[
 \Patt_{s_\tau}(F)\le[n+1+nC_\tau]m^{\tau n},\qquad |F|=n.
\]
There is no factor depending on the last observation time or on
the unbounded absorption time.

For the matching lower count prescribe any $n$ core blocks at the
positions of $F$, fill all other one-step digits with zero, and let
$x=\sum_{j\ge0}e_jm^{-j-1}<1$. There are only finitely many
nonzero digits, so each remaining base-$m$ tail is smaller than one.
Induction gives $\lfloor mT^jx\rfloor=e_j$ and shows that the
rule generates the prescribed digits. Hence $\Patt_{s_\tau}(F)\ge m^{\tau n}$ for every
window. The all-window bounds and the full-class consecutive lower
bound prove the proposition, with minima actually attained.
Control-label projection treats arbitrary atom refinements.
\end{proof}

When $m\ge3$ and $D=m$, the minimum one-step safe subcover has
$r=m+1$, although $L_1=M_1=\log m$. Indeed, $m$ intervals of
length $R/m$ could cover an interval of length $R$ only by abutting
without positive-length overlap. Their left endpoints are multiples
of $1/m$, so this would force their consecutive control differences
to equal the noninteger $R=m/(m-1)$. All $m+1$ intervals are
therefore needed. Although the legal branches have different
successors, they share the state factor $x\mapsto e^{2\pi ix}$,
with induced map $z\mapsto z^m$.

\subsection{Protected inverse branches and common sampling}
On $Q=[0,1]$, at period one, use
\[
 g_a(x)=\frac54x,\quad g_b(x)=\frac{20}3x-\frac{16}3,
 \quad g_c(x)=20x-19,\quad g_d(x)=100x.
\]
Their safety domains are
$[0,4/5]$, $[4/5,19/20]$, $[19/20,1]$, and
$[0,1/100]$, respectively. Continuous ambient extensions are
obtained by truncating these affine maps to $[-1,2]$, with no
change on legal paths. The minimum subcover is three and the largest
two inverse slopes sum to $4/5+3/20=19/20<1$; hence
$M_1=\log3$ by the scalar special case of the capacity theorem.

\begin{proposition}\label{prop:bridge}
For $A_2=(0,2,4,\ldots)$, every legal Borel feedback and every
finite Borel atom refinement has at least $3^n$ actual patterns
on its first $n$ observations. Consequently
\[
 \min_s h_{A_2}(s)=L_1=M_1=\log3.
\]
\end{proposition}
\begin{proof}
We construct two-step inverse paths whose first controls can be
chosen freely and whose second control is always $b$. Each path lies
in regions where safety forces those controls, so the construction
works for every feedback on the same even sampling sequence.

The full inverse branches are
$\psi_a(y)=4y/5$, $\psi_b(y)=4/5+3y/20$,
$\psi_c(y)=19/20+y/20$. Let $J=[1/2,999/1000]$ and
$\theta_u=\psi_u\psi_b$. They satisfy
\[
 \theta_a(y)=16/25+3y/25,\quad
 \theta_b(y)=23/25+9y/400,\quad
 \theta_c(y)=99/100+3y/400.
\]
Their complete images and the intermediate image are
\[
\begin{array}{c|c|c}
 \text{image}&\text{exact endpoints}&\text{only legal control}\\\hline
 \psi_b(J)&[7/8,18997/20000]&b\\
 \theta_a(J)&[7/10,18997/25000]&a\\
 \theta_b(J)&[149/160,376991/400000]&b\\
 \theta_c(J)&[159/160,398997/400000]&c.
\end{array}
\]
Each $\theta_u(J)$ lies in $J$. Every displayed closed interval
is strictly inside the appropriate uniquely legal region; in
particular $\theta_a(J)$ avoids the extra control $d$.

Given any $u_0\cdots u_{n-1}\in\{a,b,c\}^n$, set
$x_{2n}=1/2$ and recursively
$x_{2j+1}=\psi_b(x_{2j+2})$,
$x_{2j}=\theta_{u_j}(x_{2j+2})$.
The table shows that at every even state only $u_j$ is legal
and at every odd state only $b$ is legal. Thus \emph{every}
feedback must execute the actual period-one prefix
$u_0,b,\ldots,u_{n-1},b$ from the same constructed initial state.
The endpoint remains in $Q$ and the globally legal rule continues
safely forever. There is no added phase or period-two control.
All three-symbol sampled words are therefore actual patterns.
Refinement only increases their count.

The Borel selector using $a$ below $4/5$, $b$ from $4/5$ to
$19/20$, and $c$ at and above $19/20$, is legal at all endpoints
and uses only three labels. Its upper count is $3^n$ on every
window. This matches the lower bound, including attainment, and
$L_1\le M_1$ completes the equalities. The normalization is
$n\tau=n$, not $2n$.
\end{proof}

The closed contracting branches $\theta_a,\theta_b,\theta_c$
have disjoint first-level images in $J$. Their nested images give
a compact set homeomorphic to the three-symbol full shift, and
every $T_s^2$ on that set is the same shift under this coding.
To verify this, a word's image has diameter at most
$(3/25)^n\operatorname{diam}J$; nested images intersect in one
point, different first differing letters give disjoint images, and
the inverse identities used above give the shift relation.
The union of this set with its $\psi_b$ image is protected under
every legal feedback. These sets can overlap: the all-$b$ state
$16/17$ belongs to both and uniquely uses $b$. Hence the proof
does not introduce an implicit disjoint phase register. This is
precisely a protected-subsystem mechanism, even if the rest of the
state space has feedback-dependent successors.

\subsection{A noninteger four-branch system}
\label{subsec:four-affine}
Let $Q=[0,6]$ and use, at period one,
\[
 g_{A0}(x)=\tfrac32x,\quad g_{A1}(x)=\tfrac32x-3,
 \quad g_{B0}(x)=\tfrac53x,\quad g_{B1}(x)=\tfrac53x-4.
\]
The original safety intervals are
$[0,4]$, $[2,6]$, $[0,18/5]$, $[12/5,6]$.
Truncation to the ambient interval $[-1,7]$ retains exactly these
safe transitions. All legal Borel feedbacks, arbitrary mixtures,
and finite atom refinements are allowed. Two $A$ domains cover
$Q$ and no single one does, so $r=2$. The largest scalar
inverse capacity is $2/3<1$. Thus $M_1=\log2$.
Every actual length-$N$ control prefix has its fiber inside an
interval of length at most $6(3/2)^{-N}$; these intervals cover
$Q$. Its number $D_N(s)$ is therefore at least $(3/2)^N$.
Appendix~\ref{app:finiteblock} retains the weights and actual
even--odd compatibility in a refinement of this volume argument.
This gives
\begin{equation}\label{eq:four-L}
 \log(3/2)\le L_1\le\log2.
\end{equation}

For $A_2=(0,2,\ldots)$ let
$\ell_2=\inf_s h_{A_2}(s)$, and let $\ell_2^A$ restrict to
all two-$A$ Borel feedbacks, without restricting their partitions.
The unrefined two-$A$ rules have the exact description
\[
 s_E(x)=A0\text{ on }[0,2)\cup E,
 \quad s_E(x)=A1\text{ elsewhere},\qquad E\subset[2,4]
 \text{ Borel}.
\]
Indeed the two overlap domains are exactly $[2,4]$.
Write $N_n(s)=\Patt_s(\{0,2,\ldots,2n-2\})$ for the even
pattern count. For any feedback, an actual $2n$ prefix is determined
by its even and odd sampled patterns. Odd patterns occur among the even
patterns from states $T_sx\in Q$. Hence
$D_{2n}(s)\le N_n(s)^2$ and the volume lower bound gives
$N_n(s)\ge(3/2)^n$. Consequently $\ell_2\ge\log(3/2)$.

\begin{proposition}[Density-one samplings]\label{prop:density-one}
If $A=(a_i)$ satisfies $a_n/n\to1$, then
$\inf_s h_A(s)=\log(3/2)$ over the full class.
The greedy feedback $s_g=s_\varnothing$ is an actual minimizer
for every such $A$.
\end{proposition}
\begin{proof}
The consecutive counts of the greedy rule grow like $(3/2)^n$.
To establish the sampled statement, we first bound these counts on
the entire interval, including points that have not yet reached the
invariant core. A density-one sequence omits only $o(n)$ coordinates
from its elapsed prefix, and that omission changes the logarithmic
count by $o(n)$.

Set $\beta=3/2$. On the invariant core $[0,3)$, scale $x=3t$.
The greedy dynamics becomes $B(t)=\beta t-\lfloor\beta t\rfloor$
on $[0,1)$. Every nonempty length-$n$ digit cylinder is a
left-closed right-open interval whose $B^n$ image is $[0,l)$
for some $0<l\le1$. This is proved by splitting the current
image at $1/\beta$: it has a single child if $\beta l\le1$,
and otherwise a full-image child and a nonfull-image child.
Equality creates no extra singleton at the excluded right endpoint.

Let $F_n$ be the number of full-image cylinders, with $F_0=1$.
The sum of image lengths is $\beta^n$, since each source length
is $l/\beta^n$ and the sources partition $[0,1)$. Thus
$F_n\le\beta^n$. Each nonfull cylinder has a last full-image
ancestor, and each ancestor has at most one descendant at a given
later level without another full-image step. The total core count
$C_n$ consequently satisfies
\[
 C_n\le\sum_{j=0}^n F_j\le3\beta^n.
\]
For $x\in[3,6)$ the greedy rule uses $A1$ until
entering the core; while waiting its state is
$6-\beta^j(6-x)$, so entry is finite. State $6$ uses $A1$
forever. Therefore the full-space count satisfies
\[
 D_n(s_g)\le1+\sum_{j=1}^n C_j\le10\beta^n.
\]
This includes all thresholds and the exceptional endpoint.

Let $N=a_n+1$, so $N-n=o(n)$. Each sampled word has at most
$4^{N-n}$ completions to a consecutive word, giving
\[
 \Patt_s(\{a_1,\ldots,a_n\})\ge
 \frac{D_N(s)}{4^{N-n}}\ge\frac{\beta^N}{4^{N-n}}.
\]
For the greedy rule the same count is at most
$D_N(s_g)\le10\beta^N$. After division by $n$ the two
rates tend to $\log\beta$. Refinements preserve the lower
bound and the unrefined rule gives the matching upper bound.
\end{proof}

Let $\rho>1$ be the unique solution
\begin{equation}\label{eq:rho}
 \rho^6=\rho^5+\rho^4+\rho+1,
 \qquad\log\rho=0.55600938749\ldots.
\end{equation}
Appendix~\ref{app:greedy} computes the greedy feedback's exact rate
$h_{A_2}(s_g)=\log\rho$, with full upper and lower counts.
Combining these estimates gives
\begin{equation}\label{eq:four-even}
 \log(3/2)\le\ell_2\le\ell_2^A\le\log\rho.
\end{equation}
Thus we know the greedy rate at this one sequence, but not whether
it minimizes that rate among two-$A$ rules or among all four-control
rules. Optimizing the sequence as well gives the separate bounds
\eqref{eq:four-L}. The consistency obstruction below explains why an
open-loop language with favorable counts is not enough to improve
either upper bound.


\subsection{An alternating open-loop language with no Borel realization}
Let $R>0$, $1<a,b<2$, and $\log a/\log b\notin\mathbb Q$.
On $[0,R]$ take four legal affine branches
\[
 g_{Ai}(x)=ax-i(a-1)R,\qquad
 g_{Bi}(x)=bx-i(b-1)R\quad(i=0,1).
\]
Each type's two full inverse intervals cover $Q$. Let
$\Lambda_{\rm alt}$ be the two-phase language whose $A,B$ types
alternate strictly and whose digits are otherwise arbitrary. Let
\[
 \widehat\Lambda=\Lambda_{\rm alt}\ \cup\
 \{w c^\infty:w\text{ is a finite alternating-type word},\ c\in U\}.
\]
The empty prefix is allowed. This compact shift-invariant language
includes all constant endpoint tails and their finite alternating
prehistories. The inverse contractions
$\psi_{Ci}(y)=(y+i(\lambda_C-1)R)/\lambda_C$, with
$\lambda_A=a,\lambda_B=b$, define a continuous coding
$\pi(d)=\lim_n\psi_{d_0}\cdots\psi_{d_{n-1}}(0)$.
Every state has a coding in $\Lambda_{\rm alt}$: in each prescribed
type choose a safe digit and repeat. The contraction error is at
most $R\min(a,b)^{-n}$. This proves full open-loop coverage, not
feedback consistency.

\begin{proposition}\label{prop:alternation}
No globally safe Borel period-one feedback has all its control names
in $\widehat\Lambda$. This remains true with arbitrary Borel
mixtures and finite atom refinements. Nevertheless for every
$n$-position window $F$ with $n\ge1$,
\[
 2^{n+1}\le\Patt_{\widehat\Lambda}(F)\le8(n+1)2^n.
\]
\end{proposition}
\begin{proof}
The alternating part has two distinct type patterns and $2^n$
independent digits. For a constant-tail member there are at most
$n+1$ cuts on $F$, two prefix phases, $2^j$ prefix digits, and
four constant tails. Summing by cuts bounds the count by
$8(n+1)2^n$ (and also covers the infinite alternating case).

To prove nonrealizability, examine small positive states. There only
zero-digit controls are safe, so two alternating steps multiply the
state by $ab$. In logarithmic coordinates, a Borel choice of the first
type would give a two-coloring of a circle exchanged by an irrational
rotation. The following argument rules out that coloring.

Let $C$ be the union of all finite affine
preimages of $\{0,R\}$. This set is countable, and its complement is preserved by every
legal branch: a successor in $C$ would put the original state in $C$. A constant safe tail under any one
branch forces its tail state to be the corresponding endpoint,
by expansion about $0$ or $R$. Thus any hypothetical feedback
has strictly alternating types forever on $Q\setminus C$.
Choose
\[
 0<\eta<\min\{R/a,R/b,(a-1)R/a,(b-1)R/b\}.
\]
On $(0,\eta)$ only zero-digit controls are legal. For
$0<x<\eta/(ab)$ outside $C$, the first two steps are
$x,ax,abx$ if the first type is $A$, and $x,bx,abx$ otherwise.
The third type is the first type. Put $p=\log a,q=\log b,h=p+q$
and encode the feedback type at $\eta e^t$, $t<0$, by a Borel
$c(t)\in\{0,1\}$, with $A$ represented by one. Outside a
countable exceptional set,
\[
 c(t+h)=c(t),\quad
 c(t)=1\Rightarrow c(t+p)=0,\quad
 c(t)=0\Rightarrow c(t+q)=1\qquad(t<-h).
\]
Translating each negative interval of length $h$ to $[-h,0)$
gives a measurable circle function, with all countably many
translated exceptions still null. Apply the last two implications
at $t=r-h$, $-h\le r<0$. For its set $E$ of type one, circle
rotations satisfy $R_pE\subset E^c$ and $R_qE^c\subset E$
modulo null sets. Since $q=-p$ modulo $h$, these imply
$R_pE=E^c$, measure $1/2$, and $R_{2p}E=E$ modulo null sets.
The number $2p/h$ is irrational. For
$G=2\one_E-1$, invariance under that rotation and Fourier
coefficients give
$\widehat G(j)=e^{2\pi i j(2p/h)}\widehat G(j)$.
Every nonzero coefficient vanishes; the mean vanishes as well.
Completeness of the trigonometric system gives $G=0$ almost
everywhere, contradicting $|G|=1$. This proves impossibility.
Refinements project to the same physical control rule and cannot
avoid the obstruction.
\end{proof}

The noninteger system in Section~\ref{subsec:four-affine} has
$a=3/2,b=5/3,R=6$. Its log ratio is irrational by unique prime
factorization: $(3/2)^j=(5/3)^k$ would imply
$3^{j+k}=2^j5^k$. None of its nonempty finite affine composites
has an integer slope. Thus its alternating $A/B$ open-loop language
cannot be realized by a globally safe Borel feedback. Other feedbacks
remain available, so this obstruction does not determine the optimized
sampling rates.

\section{The remaining common-sampling problem}\label{app:scope}
The unresolved step in the capacity approach is the choice of times.
The product estimate forces independent coordinates for each
feedback, but those coordinates may move when the rule changes.
Does it nevertheless imply $L_\tau=M_\tau$ for the full Borel
feedback class, or can the two values differ?

The four-branch affine system makes this question concrete:
\[
 \log(3/2)\le L_1\le M_1=\log2.
\]
At density-one samplings the inner minimum is $\log(3/2)$; at even
times the greedy rate is $\log\rho$, with $\rho$ given by
\eqref{eq:rho}. Neither the optimal even-time rate nor $L_1$ is
determined. The weighted pattern-pair identity in
Appendix~\ref{app:finiteblock} identifies possible improvements to
the volume lower bound, but gives no uniform positive improvement.
A resolution needs additional information about compatible feedback
fibers, beyond an open-loop covering count.

\appendix

\section{Attainment of the fixed-period infimum}\label{app:minimum}

Integer-logarithm quantization of maximal-pattern entropy is a known
combinatorial fact. For one-sided subshifts it is stated in
Huang and Ye \cite[Theorem~5.5]{HuangYe2009}. We record the finite-set
argument needed here, since it also settles attainment of the feedback
infimum. This argument does not address the common-sampling problem.

\begin{proposition}\label{prop:quantization}
Let $\varnothing\ne\Omega\subset E^{\Nzero}$, where $|E|=q<\infty$,
and fix $\tau>0$. There is an integer $m\in\{1,\ldots,q\}$ such that
$\hpat(\Omega;\tau)=\log m/\tau$.
\end{proposition}
\begin{proof}
For $H\subset E^n$, let $d_k(H)$ be the largest size of a coordinate
set $J$ for which $H|_J$ contains $\prod_{j\in J}V_j$, with
$V_j\subset E$ and $|V_j|=k$. Put $d_k(\varnothing)=-1$.
The elementary finite bound is
\[
 |H|\le B_k(n,d):=
 \sum_{i=0}^{\min(d,n)}\binom ni(k-1)^{n-i}\binom qk^i
 \quad\text{if }d_k(H)\le d,\qquad 2\le k\le q.
\]
Set $B_k(n,-1)=0$. To prove the bound, write $H_e$ for the
$(n-1)$-letter prefixes ending in $e$, $V=\bigcup_eH_e$, and
$V_S=\bigcap_{e\in S}H_e$ for $|S|=k$. Projection gives
$d_k(V)\le d$, whereas $d_k(V_S)\le d-1$: a shattered set in
$V_S$ can be extended by the last coordinate. A prefix appearing
in $t$ of the sets $H_e$ contributes
$t\le(k-1)+\binom tk$. Hence
\[
 |H|\le(k-1)|V|+\sum_{|S|=k}|V_S|
 \le(k-1)B_k(n-1,d)+\binom qk B_k(n-1,d-1)
 =B_k(n,d).
\]
The empty class and the nonempty class at $n=0$ start the induction.

Let $D_k=\sup_Fd_k(\Omega|_F)$, over finite $F\subset\Nzero$,
and let $m$ be the largest $k$ with $D_k=\infty$.
Since $\Omega$ is nonempty, $D_1=\infty$.
Restricting shattered sets gives $p_\Omega^*(n)\ge m^n$ for every
$n$. If $m=q$, the alphabet bound gives equality. Otherwise
$D_{m+1}=d<\infty$, so the preceding bound gives, for a constant $C$,
\[
 m^n\le p_\Omega^*(n)
 \le\sum_{i=0}^{\min(d,n)}\binom ni m^{n-i}\binom q{m+1}^i
 \le C n^d m^n\qquad(n\ge1).
\]
Taking logarithms and dividing by $n\tau$ proves the assertion.
\end{proof}

\begin{corollary}\label{cor:pattern-minimum}
For every nonempty class of finite Borel invariant partitions at
one fixed period $\tau$, the infimum of $\hpat(\C;Q)$ is a minimum.
At the common attaining schedule of Corollary~\ref{thm:B}, the inner
sequence-entropy infimum is also a minimum. In particular, for a
finite control alphabet the quantity $M_\tau$ in
\eqref{cc-eq:LM} is attained without a countability assumption.
\end{corollary}
\begin{proof}
Apply Proposition~\ref{prop:quantization} to each $\Sigma_\C$.
The corresponding integers form a nonempty subset of $\N$, which
has a least member. At the schedule of Corollary~\ref{thm:B}, all
individual sequence rates equal their pattern rates.
For the full finite-control class, Lemma~\ref{cc-lem:merge}
identifies the infimum with that over legal selectors, whose
nonempty fibers have at most $|U|^\tau$ labels.
\end{proof}

No attainment assertion for an arbitrary fixed sequence $A$ follows
from this corollary. The integer quantization applies to maximal-pattern
entropy, and fixed-sequence entropy need not have that form.

\section{A quantitative single-step estimate}\label{app:weighted}
Assume the notation of Section~\ref{sec:capacity} and
$\Theta_k<1$. The proof below gives an explicit lower density of
independent coordinates without the omitted-letter covering lemma. For nonempty $F\subset W^n$, let $d_k(F)$ be the largest
$|J|$ for which there are sets $V_j\subset W$, $|V_j|=k$, such that
$F|_J\supset\prod_{j\in J}V_j$. The sets $V_j$ may differ with $j$.
Set $d_k(\varnothing)=-1$. Put $\theta=\Theta_k<1$, $C_0=\|\sum_u\rho_u\|_\infty$ and
$B=C_0\binom qk$. For nonnegative $f$ let
$\Phi f=\max_{|S|=k-1}\sum_{u\in S}L_uf$ and
$\mathcal Tf=\sum_uL_uf$. Their norms are bounded by $\theta$
and $C_0$, respectively.

\begin{proposition}[Weighted shattering estimate]\label{prop:weighted-sauer}
With $H(n,-1)=0$, for every $F\subset W^n$,
\begin{equation}\label{eq:cap-weighted}
 \opnorm{P_F}\le
 H(n,d_k(F)),\qquad
 H(n,d)=\sum_{j=0}^{\min(d,n)}\binom nj B^j\theta^{n-j}.
\end{equation}
For actual names this gives $d_k(F_n(s))>\alpha n$, uniformly in
every legal feedback, whenever $0<\alpha<1/2$ satisfies
\begin{equation}\label{eq:cap-alpha}
 \log\theta+\mathfrak h(\alpha)+\alpha\log(B/\theta)<0,
 \qquad\mathfrak h(t)=-t\log t-(1-t)\log(1-t).
\end{equation}
\end{proposition}
\begin{proof}
Take $H(n,-1)=0$ and treat the empty word at $n=0$ separately.
Split $F$ at its \emph{first} letter: $F_a=\{v:av\in F\}$,
$U_0=\bigcup_aF_a$, and $I_S=\bigcap_{a\in S}F_a$ for
$|S|=k$. If $d=d_k(F)$, then $d_k(U_0)\le d$ and
$d_k(I_S)\le d-1$: independence of $d$ tail positions in the latter
would also make the first position independent and give $d+1$.
For each $v\in U_0$ the set of its permitted first letters has
either at most $k-1$ members, or contains a $k$-set $S$. Applying
the positive tail product $L_v$ and summing gives
\[
 P_F f\le P_{U_0}\Phi f+\sum_{|S|=k}P_{I_S}\mathcal T f.
\]
This respects the order $L_{av}=L_vL_a$; no operators are
commuted. Induction and Pascal's identity give
$\theta H(n-1,d)+BH(n-1,d-1)=H(n,d)$. Positive domination
handles signed $f$ and proves \eqref{eq:cap-weighted}.

By actual-fiber mass, $1\le H(n,d_k(F_n(s)))$.
The cover property gives $C_0\ge1$, and positivity gives
$0<\theta\le C_0\le B$. For $d\le\alpha n$, therefore,
$B/\theta\ge1$ and
$\sum_{j\le\alpha n}\binom nj\le e^{n\mathfrak h(\alpha)}$
give
$H(n,d)\le\exp(n[\log\theta+\mathfrak h(\alpha)+
\alpha\log(B/\theta)])<1$, a contradiction.
The binomial bound follows by comparing each term to its probability
under independent Bernoulli$(\alpha)$ trials; those probabilities
sum to at most one. Condition~\eqref{eq:cap-alpha} is possible
because its last two terms tend to zero as $\alpha$ tends to zero.
\end{proof}

\subsection{A critical boundary of the single-step density condition}
For $3/2<\lambda\le2$, on $Q=[0,1]$ use
\[
 g_a(x)=\lambda x,\quad
 g_b(x)=\lambda x-(\lambda-1)/2,\quad
 g_c(x)=3x-2.
\]
The safety intervals are
\[
 K_a=[0,1/\lambda],\quad
 K_b=[(\lambda-1)/(2\lambda),(\lambda+1)/(2\lambda)],
 \quad K_c=[2/3,1].
\]
They cover $Q$ and $r=3$. Both $a$ and $c$ are forced into
every subcover by endpoints $0,1$. Since $1/\lambda<2/3$,
they leave an open gap, so $b$ is also needed. The inverse
densities are constants $1/\lambda,1/\lambda,1/3$; therefore
$\Theta_2=1/\lambda<1$ and $\Theta_3=2/\lambda$.

\begin{proposition}\label{prop:critical}
For every $3/2<\lambda\le2$, this full Borel class has
$M_1=\log2<\log3$. Thus the condition $\Theta_3<1$ cannot be
uniformly replaced by $\Theta_3\le1$ or by
$\Theta_3<1+\eta$ for any $\eta>0$, while retaining the
conclusion $M_1=\log3$.
\end{proposition}
\begin{proof}
The $k=2$ single-step criterion gives every feedback
$h^*\ge\log2$. Put $t=(\lambda-1)/(2\lambda)$ and define
\[
 s_\lambda(x)=\begin{cases}
 a,&0\le x<t,\\ b,&t\le x<2/3,\\ c,&2/3\le x\le1.
 \end{cases}
\]
It is Borel and legal: $t<1/\lambda$ and
$(\lambda+1)/(2\lambda)\ge3/4>2/3$.
For $x<t$, $g_a(x)<(\lambda-1)/2\le1/2$.
For $t\le x\le1/2$,
$g_b(x)=\lambda x-(\lambda-1)/2\in[0,1/2]$.
Hence $J=[0,1/2]$ is invariant. Before entering $J$ only
$b,c$ occur; once in $J$ only $a,b$ occur. This remains true
for states which never enter and at all thresholds, including the
fixed points $1/2$ and $1$.

At any $n$ selected times a sampled name has some cut $0\le j\le n$
between the two groups, each of size two. Thus its pattern count
is at most $(n+1)2^n$ for \emph{every} window. This actual
feedback gives $h^*\le\log2$ and completes the equality.
At $\lambda=2$, $\Theta_3=1$; as $\lambda$ increases to $2$
from below, $\Theta_3$ decreases to $1$ from above. These values
give both claimed failures of uniform relaxation. The example
does not make $\Theta_3<1$ a necessary condition for every
system, and does not determine $L_1$ for this family.
\end{proof}

\section{Finite-cylinder representation and a contraction criterion}\label{sec:cylinder}

The matrix calculation in Section~\ref{subsec:product-example} works
because all iterates of the constant function lie in a finite-dimensional
space. We give that reduction for cylinder safety domains, then one
criterion that avoids checking products individually. A common positive
left weight bounds every list product by the same geometric sequence.
This is the usual linear copositive Lyapunov argument for positive
matrices; see \cite{Debauche2022} for the broader framework. The
condition below specifies when the safety relations supply such a weight.

\subsection{An exact finite representation}
Let $D$ be a finite alphabet of size $d\ge2$, $Q=D^{\Nzero}$,
and $\mu$ its fair Bernoulli probability. The finite controls have
legal maps $g_u=P_u\sigma$, where $p_u$ permutes $D$ and $P_u$
acts in every coordinate. Suppose
\[
 K_u=\bigcup_{\xi\in R_u}[\xi],\qquad R_u\subset D^h,
 \quad h\ge2,
\]
and the $K_u$ cover $Q$. Shorter cylinders can be completed to the
same length $h$. The continuous cemetery extension used above
applies to these clopen domains. The control period is always one.

For every bounded Borel function the full operator is
\begin{equation}\label{eq:cyl-operator}
 L_uf(y)=\frac1d\sum_{a\in D}
 \one_{R_u}(a,p_u^{-1}y_0,\ldots,p_u^{-1}y_{h-2})
 f(aP_u^{-1}y).
\end{equation}
Indeed, conditioned on $x_0=a$ the tail of $x$ still has law $\mu$,
and $P_u$ preserves that law. Substitution in
\eqref{eq:cap-transfer} gives this formula with the exact safety gate.

Index source and target vectors by $D^{h-1}$. For target $\eta$
and source $\xi$, set
\begin{equation}\label{eq:cyl-matrix}
 T_u(\eta,\xi)=\frac1d\#\left\{a\in D:
 \begin{array}{l}
 (a,p_u^{-1}\eta_0,\ldots,p_u^{-1}\eta_{h-2})\in R_u,\\
 (a,p_u^{-1}\eta_0,\ldots,p_u^{-1}\eta_{h-3})=\xi
 \end{array}\right\}.
\end{equation}
When $h=2$ the second tuple is just $a$.
Functions of the first $h-1$ coordinates form an invariant subspace
of every $L_u$, represented by $T_u$. Starting at $\one$, every
list product stays in this space. Every such cylinder has positive
measure and a positive product has norm equal to its value on $\one$.
Therefore, for every length and every history,
\begin{equation}\label{eq:cyl-fullnorm}
 \opnorm{A_{S_{n-1}}\cdots A_{S_0}}
 =\|T_{S_{n-1}}\cdots T_{S_0}\one\|_\infty,
 \qquad T_S=\sum_{u\in S}T_u.
\end{equation}
Only the operator calculation has become finite-dimensional; the
feedback itself may still depend on the whole state.

For two-letter domains, rows are target letters and columns are
source letters, with
\begin{equation}\label{eq:cyl-two}
 T_u(j,i)=d^{-1}\one_{R_u}(i,p_u^{-1}j),\qquad
 d(w^\top T_u)_i=\sum_{z:(i,z)\in R_u}w(p_uz).
\end{equation}
The structural conditions below use this source-column identity.
A left budget $w^\top T_S\le\theta w^\top$ is different from
changing a reference density, which leads instead to a right
inequality $T_Sa\le\theta a$. They cannot be interchanged.

\subsection{A strict source-weight criterion}
Partition $D=A\sqcup B\sqcup C$ into nonempty sets of sizes
$a,b,c$, so $d=a+b+c$. Use three controls with
\begin{equation}\label{eq:weight-land}
 p_0(C)\subset D\setminus A,\qquad p_1(B)\subset A,
 \qquad p_2(A)\subset D\setminus A.
\end{equation}
Choose arbitrary relations $E_0\subset A\times C$ and
$E_1\subset A\times B$, with $E_0\cup E_1\ne\varnothing$,
and define
\begin{align}
 R_0&=((D\setminus A)\times C)\cup E_0,\nonumber\\
 R_1&=((D\setminus A)\times B)\cup E_1,\label{eq:weight-domains}\\
 R_2&=(A\times D)\cup(D\times A),\qquad K_u=[R_u].\nonumber
\end{align}
If the source letter is in $A$, control $2$ permits every second
letter, while $E_0$ and $E_1$ specify the extra choices of controls
$0$ and $1$. Outside $A$, the second letter determines the unique
legal control. For $i\in A$ and $u=0,1$, write
\[
 e_u(i)=|\{z:(i,z)\in E_u\}|,\qquad e=\max_{i\in A}e_1(i).
\]
The following inequality will allow a single choice of weights to
control all two-label lists:
\begin{equation}\label{eq:weight-count}
 s_-:=\max\left\{\frac c{b+c},\frac e{b+c-e}\right\}
 <s_+:=\frac{\min\{a,c\}}b.
\end{equation}
Its lower bound limits the extra choices on $A$; its upper bound
limits the weight transferred from sources outside $A$.

\begin{proposition}[Strict source budget]\label{thm:source-weight}
For every system satisfying \eqref{eq:weight-land}--\eqref{eq:weight-count},
$r=3$, all histories of global two-label lists have exponentially
decaying inverse capacity, and every legal unrefined Borel feedback
has $h^*(s)=M_1=\log3$. Every finite Borel atom refinement has
$h^*\ge\log3$. Moreover, no full-support reference probability
can directly satisfy the joint-density or scalar single-step
criterion for $k=3$ on these original branches.
\end{proposition}
\begin{proof}
We seek one weighted sum of the source coordinates that contracts
under every list. Giving the class $A$ a larger weight compensates
for its extra legal edges. The interval $(s_-,s_+)$ is precisely
what leaves room to choose that weight.

Let $s>0$ be chosen below, put $\gamma=1+s$, and give weight
$w_i=\gamma$ to $i\in A$ and weight $w_i=1$ to the other
letters. The weighted mass of the constant vector is
$H=\sum_iw_i=d+sa$.
Before division by $d$, the weighted contributions of controls
$0,1,2$ in a source column are
\[
\begin{array}{c|ccc}
 &0&1&2\\\hline
 i\in A&e_0(i)\le c&\gamma e_1(i)\le\gamma e&H\\
 i\notin A&c&\gamma b&a.
\end{array}
\]
Each entry follows from the original domains and
\eqref{eq:weight-land}. For example, on $A$ control $2$ permits
every tail, hence the sum of all weights irrespective of its
permutation. Off $A$, its permitted tails lie in $A$ and land
outside $A$, each with weight one.

For a source outside $A$, the two nontrivial requirements that
each pair in the second row have sum less than $d$ are
\[
 c+\gamma b<d,\qquad \gamma b+a<d,
 \quad\text{or equivalently}\quad sb<\min\{a,c\}.
\]
For a source in $A$, compare a pair containing the contribution
$H$ with the available weight $d\gamma$. Since
$d\gamma-H=s(b+c)$, it is enough to require
\[
 c<s(b+c),\qquad \gamma e<s(b+c).
\]
These are exactly $s>s_-$, while the preceding upper restriction
is $s<s_+$. Condition~\eqref{eq:weight-count} therefore permits
$s=(s_-+s_+)/2$. With this choice $c<H$ and
$\gamma e\le\gamma b<H$, so the pair excluding $H$ is also
bounded by $H+\max\{c,\gamma e\}$. Define
\[
 \theta=\max\left\{
 \frac{H+\max\{c,\gamma e\}}{d\gamma},
 \frac{c+\gamma b}{d},\frac{c+a}{d},\frac{\gamma b+a}{d}
 \right\}.
\]
The inequalities just established, together with $a+c<d$,
give $0<\theta<1$ and
$w^\top T_S\le\theta w^\top$ for every global two-label list.
Repeatedly applying this single inequality gives
\[
 \|T_{S_{n-1}}\cdots T_{S_0}\one\|_\infty
 \le w^\top T_{S_{n-1}}\cdots T_{S_0}\one
 \le H\theta^n,
\]
since $\min_iw_i=1$. Equation~\eqref{eq:cyl-fullnorm} controls
the complete capacity for all lengths.

If $x_0\in A$, control $2$ is legal; otherwise its next letter
in $A,B,C$ requires, respectively, $2,1,0$. Thus the domains
cover $Q$. Taking $x_0\notin A$ and $x_1$ successively in
$C,B,A$ gives cylinders on which each of $0,1,2$ alone is legal,
so $r=3$. A clopen attaining selector chooses $2$ when $x_0\in A$
or $x_1\in A$, otherwise $1$ when $x_1\in B$, otherwise $0$.
It is safe, single-valued, and stays in $Q$ at each step; the same
rule at every successor gives all-time safety. Applying
Theorem~\ref{thm:capacity-main} proves all pattern conclusions,
including the lower bound for arbitrary refinements.

Let $O=[E_0\cup E_1]$. It is nonempty and clopen, and safety
multiplicity is exactly $1+\one_O$. For a full-support probability
$\nu$, either some branch fails the bounded-density prerequisite,
or all densities exist. In the latter case the images of controls
$0$ and $1$ lie, respectively, in $[D\setminus A]$ and $[A]$;
their supports are disjoint. The largest two densities consequently
sum to all three almost everywhere. Their integral is
$\sum_u\nu(K_u)=1+\nu(O)>1$. Thus $\Theta_3(\nu)>1$.
The scalar condition implies $\Theta_3<1$ and also fails.
\end{proof}

For instance
$1\le b\le a\le b+c$ and $c>b$ imply
\eqref{eq:weight-count} for every possible $e\le b$:
$e/(b+c-e)\le b/c<1$, $c/(b+c)<1$, while $s_+\ge1$.
Every allowed overlapping relation and permutation landing pattern
is covered. Constants can be uniform when these counts have fixed
bounds.

For a concrete example take
$D=\{0,1,2,3,4\}$, $A=\{0,1\}$, $B=\{2\}$,
$C=\{3,4\}$, $E_0=\{(0,3),(1,4)\}$,
$E_1=\{(0,2)\}$. Let $p_0$ be the identity,
$p_1=(0\ 2)$, and specify $p_2$ by
$0\mapsto2$, $1\mapsto3$, $2\mapsto0$, $3\mapsto4$,
$4\mapsto1$. Here $s=4/3$, $\gamma=7/3$, $H=23/3$,
$\theta=13/15$, and
$b_n\le(23/3)(13/15)^n$. The target densities at first letters
$0,1,2,3,4$ are, respectively,
\[
 (0,4/5,2/5),\quad(0,0,2/5),\quad(0,0,1),\quad
 (4/5,0,1),\quad(4/5,0,2/5).
\]
Thus $\Theta_3=9/5$. In $[14]$ the two legal successors have
first letters $4$ and $1$. No common successor is present. The
natural $A/B/C$ merging also fails to commute with the listed
permutations; this observation does not exclude every possible factor.

\section{Exact even-sampling entropy of the greedy feedback}
\label{app:greedy}

We prove $h_{A_2}(s_g)=\log\rho$ for the system of
Section~\ref{subsec:four-affine}. The proof counts words in two
directions. Certain forbidden patterns bound every greedy name from
above. Parsing the remaining words into three blocks produces a
recurrence with growth rate $\rho$. For the lower bound we realize
complete block strings by inverse paths of the same greedy map;
we do not need every word allowed by the upper bound to occur.
On $C_0=[0,3)$ the actual map $R=T_{s_g}^2$ is
\begin{equation}\label{eq:greedy-R}
 R(x)=\begin{cases}
 9x/4,&0\le x<4/3,\\
 9x/4-3,&4/3\le x<2,\\
 9x/4-9/2,&2\le x<3.
 \end{cases}
\end{equation}
Each branch is a composition of two steps of $s_g$: at the
intermediate state the same threshold $2$ decides between $A0$
and $A1$. Let $0,1$ denote the even-time labels $A0,A1$.

If a core even name begins $11$, its second state is in
$[2,9/4)$, its third in $[0,9/16)$, and the next two images
are in $[0,81/64)$ and $[0,729/256)$. The first of these
is below $4/3$. If the last image is at least $2$, its next
image is smaller than $1953/1024<2$.
Consequently every core even name avoids
\begin{equation}\label{eq:greedy-forbidden}
 111,\qquad1101,\qquad110011.
\end{equation}
All interval upper endpoints in these calculations are excluded by
the core or its actual branch convention; $4/3$ and $2$ both
map to zero under $R$.

Every word avoiding \eqref{eq:greedy-forbidden} has a unique
left-to-right parsing into the prefix-code blocks
$a=0$, $b=10$, $c=1100$, with no adjacent $cc$, followed by
one of $\varnothing,1,11,110$. The first two exclusions bound
one-runs and force two zeros after a full double-one run; the third
prevents two complete $c$ blocks. Let $H_m$ count complete block
strings of bit length $m$ with no $cc$, with $H_m=0$ for
$m<0$. If $X$ and $Z$ count strings ending in $a/b$ and in
$c$, respectively, their formal generating series satisfy
\[
 X=(z+z^2)H,\qquad Z=z^4(1+X),\qquad H=1+X+Z.
\]
Thus
\begin{equation}\label{eq:greedy-H}
 H(z)=\frac{1+z^4}{1-z-z^2-z^5-z^6},\qquad
 H_m=H_{m-1}+H_{m-2}+H_{m-5}+H_{m-6}\quad(m\ge5),
\end{equation}
with $H_0,\ldots,H_4=1,1,2,3,6$ and $H_5=10$.
The equation for $\rho$ in \eqref{eq:rho} is equivalent to
$\rho^{-1}+\rho^{-2}+\rho^{-5}+\rho^{-6}=1$, whose left side
is strictly decreasing. Its unique root obeys $5/3<\rho<7/4$.
The initial counts are at most $\rho^m$, and the recurrence
proves $H_m\le\rho^m$ at all lengths. The incomplete-tail
count is at most
$H_m+H_{m-1}+H_{m-2}+H_{m-3}\le D\rho^m$, where
$D=1+\rho^{-1}+\rho^{-2}+\rho^{-3}$.

For a real full-space even word, either all labels are $1$, or its
first $0$ occurs at some position $j$. At that moment the state
is below $2$ and lies in the core; its remaining word is governed
by \eqref{eq:greedy-forbidden}. This includes arbitrarily long
waiting outside the core and the fixed endpoint $6$.
Therefore
\begin{equation}\label{eq:greedy-upper}
 N_n(s_g)\le\sum_{m=0}^n D\rho^m
 \le\frac{D\rho}{\rho-1}\rho^n.
\end{equation}
\subsection*{Inverse paths furnishing the matching lower bound}
The three blocks are $a=0$, $b=10$ and $c=1100$. To realize their
concatenations, the set of admissible endpoints must change after a
$c$ block: a second $c$ is forbidden. We use two sets $C,D_0$ to
keep track of which inverse path is available. They are subsets of
the original state interval, not extra coordinates of the feedback.

Put $\eta=4/9$, $L=162/65$, $d_0=729/256$, and define
\[
 P_0=[0,16/27]\cup[72/65,4/3),\qquad P_1=[L,3),
\]
\[
 C=P_0\cup P_1,\qquad
 D_0=P_0\cup[2,550/243]\cup[L,d_0).
\]
Here $\eta$ is the inverse slope of $R$. The endpoint $L$ satisfies
$L=2+\eta^2L$, so the $1000$ inverse path preserves the lower
endpoint of $P_1$. Applying the $a$ inverse to $P_1$ gives
$[\eta L,4/3)=[72/65,4/3)$, explaining the second component
of $P_0$. The upper cutoff $d_0$ is chosen so that
$26/9+\eta^4d_0=3$: the $c$ inverse must stay below the excluded
endpoint of the core. The checks below verify all other endpoints.

The table gives a source for each target $y$. Its last column records
all one-step controls, including those between the observed even times.
For $b=10$ there are two inverse paths, with control words $1000$
and $1001$. The first works for every target in $C$, but its source
need not lie in $C$. When the source is required to lie there, we use
$1000$ for targets in $P_1$ and $1001$ for targets in $P_0$.
\begin{center}
\begin{tabular}{cclc}
\toprule
Even block&Target $y$&Source state&One-step controls\\
\midrule
$a=0$&$y\in C$&$\eta y$&$00$\\
$b=10$&$y\in C$&$2+\eta^2y$&$1000$\\
$b=10$&$y\in P_0$&$70/27+\eta^2y$&$1001$\\
$c=1100$&$y\in D_0$&$26/9+\eta^4y$&$10100000$\\
\bottomrule
\end{tabular}
\end{center}
The $1000$ path has its source in $D_0$ for every $y\in C$,
and in $C\cap D_0$ when $y\in P_1$. The $1001$ path supplies
a source in $C$ for every target in $P_0$.

We check the path and its source set in each row. For $a$, the source
$\eta y$ is in $P_0$, the intermediate state is $2y/3<2$,
and the endpoint is $y$; both controls are zero.
For the $1000$ path, the states after its source are
$8y/27$, $\eta y$, $2y/3$, $y$. If $y\in P_0$, the
source is in $[2,550/243)$; if $y\in P_1$, it is in
$[L,70/27)$. Each source is at least $2$, while the three subsequent control states are
below $2$. This gives $1000$ with the stated source memberships.

For the $1001$ path the states are
$8/9+8y/27$, $4/3+\eta y$, $2+2y/3$, $y$.
The source belongs to $[70/27,694/243)\subset P_1$.
The first two subsequent control states are below $2$.
The last control state $2+2y/3$ lies in $[2,26/9)$.
Thus it uses control $1$, proving $1001$ at all included endpoints.

For $c$, its source is in $[26/9,3)\subset P_1$, since
$26/9+\eta^4d_0=3$ and $d_0$ is excluded from $D_0$.
Its successive states are
\[
\begin{array}{llll}
4/3+(3/2)\eta^4y<3/2,&2+\eta^3y<9/4,&
(3/2)\eta^3y<3/8,&\eta^2y<9/16,\\
(3/2)\eta^2y<27/32,&\eta y<81/64,&
(3/2)\eta y<243/128<2,&y.
\end{array}
\]
Only its source and its second subsequent state are at least $2$;
this proves $10100000$ with every intermediate state safe.
The table now gives the relations
\[
 C\xrightarrow{a,b}C,\qquad C\xrightarrow cD_0,
 \qquad D_0\xrightarrow{a,b}C.
\]
Here an arrow $S\xrightarrow{w}T$ means that \emph{every}
endpoint in $T$ has a source in $S$ whose greedy path follows
block $w$. It is this inverse coverage of the target, not an
arbitrary forward choice at each source, that lets us concatenate
the paths.

Given a complete block string without $cc$, label its initial
boundary $C$, each boundary after $c$ by $D_0$, and each boundary
after $a$ or $b$ by $C$. All successive labels then follow one of
the displayed arrows. Choose the final state to be $0$, which
belongs to both sets, and construct the earlier states backwards.
For a $b$ step starting in $D_0$, use the $1000$ path.
For one starting in $C$, use $1000$ on target set $P_1$
and $1001$ on $P_0$. The $a$ and $c$ inverses are already
specified in the table. Each backward step lands in the set
required by the preceding arrow.

Running the greedy rule forwards from the resulting initial state
therefore produces the prescribed even word, reaches zero, and
continues safely forever. Distinct complete block strings give
distinct even words because $\{0,10,1100\}$ is a prefix code.
Hence $N_m(s_g)\ge H_m$.

Let $c_* =\min_{0\le j\le5}H_j/\rho^j>0$. The full recurrence
for $m\ge6$ proves $H_m\ge c_*\rho^m$. This and
\eqref{eq:greedy-upper} yield
\begin{equation}\label{eq:greedy-exact}
 h_{A_2}(s_g)=\log\rho.
\end{equation}
\section{Arithmetic sampling and finite-block optimization}
\label{app:finiteblock}
For a fixed globally safe Borel rule and arithmetic sampling
$(0,p,2p,\ldots)$, a length-$n+m$ sampled pattern is determined
by its first $n$ entries and the length-$m$ pattern starting from
$T_s^{pn}x\in Q$. Consequently $N_{n+m}(s)\le N_n(s)N_m(s)$.
The logarithmic rate is a limit and equals
$\inf_n n^{-1}\log N_n(s)$ (with division by $\tau$ when
the control period is $\tau$). This uses forward safety, not
surjectivity or unrestricted concatenation.

For the full four-control class at $A_2$, let $K_n$ be the least
actual count $N_n(s)$ among unrefined legal Borel feedbacks.
It exists as a least member of a nonempty subset of
$\{1,\ldots,4^n\}$. Merging refinements cannot increase counts,
and unrefined rules are themselves allowed, so
\begin{equation}\label{eq:finiteblock}
 \ell_2=\inf_s\inf_n\frac{\log N_n(s)}n
 =\inf_n\frac{\log K_n}n.
\end{equation}
The equality exchanges two infima and yields an asymptotic upper
bound from any globally implemented finite-block feedback. The
minimizing feedback may depend on $n$, so it gives no asymptotic
minimizer. Applying the formula requires actual feedback fibers,
rather than an open-loop word table.


\subsection{Weighted actual pattern pairs}
The estimate $D_{2n}\le N_n^2$ discards two kinds of information:
some patterns use the faster $B$ branches, and some even--odd pairs
never occur on one trajectory. We retain both. The feedback in this
subsection is an unrefined Borel rule for the original four-branch
system; all counts refer to its physical labels.
Set $a=3/2$, $b=5/3$, and $r_1=a/b=9/10$. Let
\[
 V_n=\{(s(T_s^{2j}x))_{0\le j<n}:x\in Q\},\qquad N_n=|V_n|,
\]
and let $I_n$ consist of the pairs $(\alpha,\gamma)$ of even
and odd patterns from the \emph{same} initial state. Since
$T_sx\in Q$, one has $I_n\subset V_n\times V_n$; surjectivity
of $T_s$ is unnecessary. Let $B(\alpha)$ count the $B0,B1$
labels of $\alpha$, and define
\[
 W_n=\sum_{\alpha\in V_n}r_1^{B(\alpha)},\qquad
 Z_n=\sum_{(\alpha,\gamma)\in I_n}
            r_1^{B(\alpha)+B(\gamma)},\qquad
 \kappa_n=Z_n/W_n^2.
\]
A $B$ label contributes the smaller inverse slope, hence the factor
$r_1$. The ratio $\kappa_n$ measures the weight of realized pairs
relative to all pairs of realized patterns. Thus
$0<\kappa_n\le1$ and $W_n\le N_n$.

\begin{proposition}\label{prop:weighted-pairs}
With the preceding definitions, put
\[
 \theta_n=-\frac1n\log\frac{W_n}{N_n},\qquad
 d_n=-\frac1{2n}\log\kappa_n.
\]
Then $\theta_n,d_n\ge0$ and
\begin{equation}\label{eq:weighted-pairs}
 \frac{\log N_n}{n}\ge\log a+\theta_n+d_n.
\end{equation}
More precisely, if $\Lambda_w$ is the slope of the complete
affine word $w$ and
\[
 \Omega_m=\sum_{w\text{ actual},\ |w|=m}\Lambda_w^{-1},
 \qquad e_n=\frac{\log\Omega_{2n}}{2n},
\]
then $\Omega_m\ge1$, $e_n\ge0$, and
\begin{equation}\label{eq:weighted-pairs-exact}
 \frac{\log N_n}{n}-\log a=\theta_n+d_n+e_n.
\end{equation}
\end{proposition}
\begin{proof}
For each actual complete word $w$, its Borel initial fiber
$E_w$ has diameter at most $6/\Lambda_w$, because two states
following the same word separate by the factor $\Lambda_w$
and their final states lie in $[0,6]$. All these fibers partition
$Q$. Consequently
\[
 1\le\sum_{w\text{ actual},\ |w|=2n}\Lambda_w^{-1}
       =a^{-2n}Z_n.
\]
The equality uses the bijection between complete words and their
actual interlaced pairs. Nonempty zero-measure fibers are included.
Substitution of $Z_n=\kappa_nW_n^2$ proves
\eqref{eq:weighted-pairs}; the same substitution with
$\Omega_{2n}=a^{-2n}Z_n$ proves
\eqref{eq:weighted-pairs-exact}.

Each original inverse branch carries $Q$ into its safety domain.
Write $G_w=g_{w_{m-1}}\circ\cdots\circ g_{w_0}$ for the affine
composition associated with a word of length $m$.
Thus the complete inverse interval $J_w=G_w^{-1}(Q)$ is contained
in $Q$ and has length $6/\Lambda_w$; every intermediate open-loop
state is safe. The feedback fiber satisfies $E_w\subset J_w$.
The intervals $J_w$ for actual words cover $Q$, and, writing
$\nu_m=\sum_w\one_{J_w}$, one obtains
\[
 6(\Omega_m-1)=\int_Q(\nu_m-1)\,dx
              =\sum_w\bigl(|J_w|-|E_w|\bigr)\ge0.
\]
Here $|E_w|$ denotes Lebesgue measure; no interval or convexity
assumption is made on a feedback fiber. This also proves the
claimed nonnegativity at every length.
\end{proof}

In \eqref{eq:weighted-pairs-exact}, $\theta_n$ is the cost of
fast-branch weighting, $d_n$ the cost of missing pairs, and $e_n$
the excess of complete inverse coverage over unit mass. They are
support-counting quantities, rather than empirical frequencies
or Shannon entropies. The estimates transfer to arbitrary atom
refinements by projecting to their physical control labels;
the equalities themselves refer to those merged labels.
Splitting an actual trajectory at time $2n$ also gives
\[
 N_{n+m}\le N_nN_m,\qquad W_{n+m}\le W_nW_m,
 \qquad Z_{n+m}\le Z_nZ_m.
\]
The weights multiply under this split, and the suffix is realized
from a state in $Q$. The subadditivity argument at the beginning
of this appendix therefore gives finite limits
\[
 h=\lim_n\frac{\log N_n}{n},\quad
 p=\lim_n\frac{\log W_n}{n},\quad
 t=\lim_n\frac{\log Z_n}{2n},\qquad
 \log a\le t\le p\le h,
\]
and the exact limiting decomposition
$h-\log a=(h-p)+(p-t)+(t-\log a)$.

These formulas do not supply a uniformly positive correction.
Even the uniformly strict one-step inequality
$\Omega_1\ge6/5$ does not iterate: controls at $0$ and $6$
are different and each inverse slope is at least $3/5$, which
proves that inequality, but the greedy feedback has
$\Omega_1=4/3$ and
\[
 1\le\Omega_m(s_g)=a^{-m}D_m(s_g)\le10.
\]
The last bound is the all-initial-state word estimate in
Section~\ref{subsec:four-affine}. Hence $e_n(s_g)\to0$.
The complete inverse intervals overcount the true fibers at one
step without producing an exponential overcount at later steps.

\begin{example}[Full support from a countable modification]
\label{ex:periodic-support}
There is a globally safe two-$A$ Borel feedback for which
$N_n=2^n$, $I_n=V_n\times V_n$, and $\theta_n=d_n=0$ at
every length. It agrees with the greedy feedback on the full
infinite trajectories of all but countably many initial states.
\end{example}
\begin{proof}
For a purely periodic binary sequence $d$, set
\[
 \pi(d)=3\sum_{j\ge0}d_ja^{-(j+1)}.
\]
This address is injective on all purely periodic sequences.
Indeed, if two different sequences share a common period $L$
and the same address, subtracting their periodic-sum formulas
gives a nonzero integer polynomial $P$ with leading coefficient
$\pm1$ and $P(3/2)=0$. If its degree is $m$, multiplication by
$2^m$ leaves an odd leading term $\pm3^m$ and only even
remaining terms, a contradiction. Degree zero is immediate.

Let $P_i$ be the addresses of the periodic sequences with first
digit $i$, and $P=P_0\cup P_1$. These sets are countable Borel
sets, and $P_0\cap P_1=\varnothing$. Define $s_{\mathrm{per}}=Ai$
on $P_i$ and $s_{\mathrm{per}}=s_g$ on $Q\setminus P$. The identity
\[
 a\pi(d)-3d_0=\pi(\sigma d)\in P\subset Q
\]
proves safety on $P$, including $0$ and $6$; safety elsewhere
is that of $s_g$. This one state rule is therefore safe on all
infinite trajectories. Injectivity and induction show that its
actual name from $\pi(d)$ is precisely $d$.

Any prescribed binary pattern on a nonempty finite window extends to a
purely periodic sequence: choose a period larger than its largest
coordinate and fill the other digits arbitrarily. The empty window has
one pattern. Thus
$\Patt_{s_{\mathrm{per}}}(F)=2^{|F|}$ for every finite $F$.
Applying this to interlaced even and odd patterns gives
$V_n=\{A0,A1\}^n$, $I_n=V_n^2$, $W_n=2^n$, $Z_n=4^n$,
and the asserted zero corrections.

Finally $R=\bigcup_{j\ge0}T_g^{-j}P$ is countable: each point
has at most two one-step greedy preimages. Outside $R$ the
greedy trajectory never enters $P$, so induction gives identical
states and controls under the two rules at every time. Nevertheless,
their consecutive-time counting entropies are respectively
$\log(3/2)$, by Proposition~\ref{prop:density-one}, and $\log2$.
Discarding countably many initial states would therefore change
the counting question. No realization of all infinite nonperiodic
binary names is needed or claimed.
\end{proof}

This example excludes an unconditional positive lower bound for
$\theta_n+d_n$ over the whole feedback class. Its own pattern
count is already large, so it does not exclude an improved estimate
restricted to feedbacks with small $N_n$, and it supplies no
low-entropy feedback or strict full-class minimax gap.


\section{The obstruction to pointwise common attainment}
\label{app:pointwise}

Corollary~\ref{thm:B} attains the rate of every member of a countable
language family. For uncountable families, pointwise common
attainment can fail even when the full feedback class has an exact
common minimum. We first describe the pattern obstruction, then realize the relevant
distinction in a system with four controls.

\subsection{The combinatorial obstruction}

\begin{example}[A purely combinatorial uncountable obstruction]
\label{ex:uncountable}
For every infinite $B\subset\Nzero$, let
\begin{align*}
 \Omega_B=\{x\in\{0,1\}^{\Nzero}:\supp(x)\subset B\}.
\end{align*}
At scale $\tau=1$,
\begin{align*}
 \Patt_{\Omega_B}(F)=2^{\#(F\cap B)},
 \qquad \hpat(\Omega_B;1)=\log2.
\end{align*}
For $A=(a_i)\in\mathscr S$,
\begin{align*}
 \hseq^A(\Omega_B;1)
 =\log2\cdot\limsup_{n\to\infty}
 \frac{\#(B\cap\{a_1,\ldots,a_n\})}{n}.
\end{align*}
If $\Nzero\setminus\{a_i:i\ge1\}$ is infinite, choose an infinite $B$ in that
complement.  Otherwise choose $B=\{a_{k^2}:k\ge1\}$.  In either case the
displayed upper density is zero.  Therefore
\begin{equation}\label{eq:abstract-obstruction}
 \sup_{A\in\mathscr S}
 \inf_{\substack{B\subset\Nzero\\B\text{ infinite}}}
 \hseq^A(\Omega_B;1)=0
 <\log2
 =\inf_{\substack{B\subset\Nzero\\B\text{ infinite}}}
 \hpat(\Omega_B;1).
\end{equation}
This proves only that minimax exchange is not a formal consequence of
pointwise attainment for arbitrary uncountable pattern families.  Every
invariant-partition name system satisfies
$\sigma(\Sigma_\C)\subset\Sigma_\C$ by suffix closure.  In contrast,
$\sigma(\Omega_B)\subset\Omega_B$ forces $k-1\in B$ whenever
$k\in B\setminus\{0\}$; an infinite such $B$ must equal $\Nzero$.  Hence the
proper sparse sets producing \eqref{eq:abstract-obstruction} cannot occur as
these autonomous control name systems.  The equation is not a deterministic
control counterexample and yields no failure statement for the unrestricted
Borel invariant-partition class.
\end{example}

\subsection{A continuous family of sparse hereditary subshifts}

We need an uncountable family of shift-invariant languages with two
features. Each language must contain full binary patterns on its own
sparse set of positions, but a fixed sampling sequence should see few
of those positions for most parameters. We encode a parameter by a
branch of the binary tree. Powers of four place its nodes so that a
pair of positions determines their translation uniquely.

Write $\Theta=\{0,1\}^{\N}$, with its fair Bernoulli probability
$\nu$. Let $\operatorname{val}_2(v)$ be the integer represented by
the binary word $v$, with value zero for the empty word. Set
\[
 j(v)=2^{|v|}+\operatorname{val}_2(v),\qquad
 b_v=4^{j(v)},\qquad
 B_\theta=\{b_{\theta|_\ell}:\ell\ge0\}.
\]
The index $j$ lists the binary-tree nodes level by level: words of
length $\ell$ receive indices $2^\ell,\ldots,2^{\ell+1}-1$.
Thus $B_\theta$ selects one marked position from each level along
the branch $\theta$.

In $K=\{0,1\}^{\Nzero}$ define
\[
 Y_\theta=
 \operatorname{cl}\bigcup_{t\in\mathbb Z}
 \{y:\supp(y)\subset(B_\theta+t)\cap\Nzero\}.
\]
We allow any subset of each translated support. Consequently these
sets are hereditary: changing any collection of $1$'s to $0$'s
keeps a sequence in the set. The closure retains this property.

\begin{lemma}\label{lem:sparse-bundle}
Each $Y_\theta$ is a nonempty hereditary subshift with
$\sigma Y_\theta=Y_\theta$ and zero ordinary entropy.
The map $\theta\mapsto Y_\theta$ is Hausdorff continuous and injective.
Every $Y_\theta$ has coordinate maximal-pattern entropy $\log2$.
For each fixed sampling sequence $A=(a_i)$,
\[
 \frac1n\log\Patt_{Y_\theta}(\{a_1,\ldots,a_n\})
 \longrightarrow0
 \quad\text{for }\nu\text{-almost every }\theta.
\]
The exceptional set may depend on $A$.
\end{lemma}
\begin{proof}
Two marked positions will determine at most one translate of the
support set. This makes finite languages locally constant in the
parameter and limits how many positions of a fixed sampling window
can belong to a typical support. We use these two consequences in
turn.

The positive differences of $D=\{4^j:j\ge1\}$ have unique
representations. Indeed, the $4$-adic valuation of
$4^j-4^i$ with $j>i$ is $i$, which determines $i$ and then $j$.
A finite pattern appears in $Y_\theta$ exactly when its set of
one-positions is contained in a translate of $B_\theta$;
cylinders are clopen and prescribed zero-positions impose no
further restriction. If there are at least two one-positions,
their difference fixes the candidate nodes and the translation.
The remaining one-positions either give incompatible nodes, or
the appearance parameters form the cylinder determined by the
deepest node. Patterns with at most one one appear for every
parameter. Thus each finite prefix language is locally constant
in $\theta$, proving Hausdorff continuity. The pattern $11$ on
$\{0,b_v-b_\varnothing\}$ detects whether $\theta$ extends a
nonempty word $v$, proving injectivity.

Setting coordinates to zero preserves the defining supports and
then their closure. Shifting replaces $t$ by $t-1$, while
prepending zero replaces it by $t+1$; compactness supplies
predecessors of limit points. This proves heredity and surjectivity.
An interval of length $N$ meets any translate of $D$ in at most
$L_N=2+\lceil\log_4(N+1)\rceil$ points. Consequently
\[
 |\mathcal L_N(Y_\theta)|
 \le (L_N+1)N^{L_N}\qquad(N\ge2),
\]
which gives zero ordinary entropy. On any $n$ positions in
$B_\theta$ all binary patterns occur, so $p_{Y_\theta}^*(n)=2^n$.

It remains to estimate patterns on a fixed sampling sequence.
For $F_n=\{a_1,\ldots,a_n\}$ put
$M_\theta(F_n)=\sup_t|F_n\cap(B_\theta+t)|$.
This is the largest possible number of $1$'s in a pattern on $F_n$.
A bound on $M_\theta(F_n)$ will therefore bound the number of
patterns by counting binary words with at most that many $1$'s.
We first list the possible translations without fixing $\theta$. If a translate meets $F_n$ twice, the two positions
give a difference $4^j-4^i$. Its unique representation determines
both marked positions and the translation. There are therefore at
most $\binom n2$ candidate translations.

We now vary $\theta$ while keeping $F_n$ fixed. For one candidate
$t$, there are at most $n$ binary-tree nodes $v$ with
$b_v+t\in F_n$. If a branch contains at least $k$ of them,
one of these nodes has depth at least $k-1$, since a branch has
only one node at each depth. The set of parameters extending a
particular node $v$ has measure $2^{-|v|}$. A union bound, first
over the nodes and then over the candidate translations, gives
\[
 \nu\{M_\theta(F_n)\ge k\}
 \le \binom n2 n\,2^{-(k-1)}
 \le n^3 2^{-(k-1)}\qquad(k\ge2).
\]
These events are finite unions of parameter cylinders. With
$k_n=\lceil6\log_2(n+1)\rceil+2$, the probabilities are summable.
Borel--Cantelli gives $M_\theta(F_n)<k_n$ eventually almost surely.
Every occurring pattern then has at most $k_n$ one-positions, and
\[
 \Patt_{Y_\theta}(F_n)\le(k_n+1)n^{k_n}\qquad(n\ge2).
\]
Its normalized logarithm tends to zero.
\end{proof}

\subsection{Four controls and a positive full-class minimum}

\begin{proposition}\label{prop:finite-pointwise}
There is a compact system with four distinct continuous control
maps and a controlled-invariant compact set $Q$ such that
\[
 \hinv(Q)=0,\qquad L_1=M_1=\log2,
\]
yet no sampling sequence attains $\hpat(\C;Q)$ for every legal
period-one Borel invariant partition $\C$.
This already fails for a family of four-label selectors with
Hausdorff continuous, compact name languages and rate $\log4$.
\end{proposition}
\begin{proof}
The first control component will be forced by a Thue--Morse state;
this fixes the full-class minimum at $\log2$. A second component
can record one of the sparse languages above, raising individual
maximal-pattern rates without permitting one sequence to attain them
all. A register makes the two choices physically distinct controls.

Let $Z$ be the Thue--Morse subshift of
Section~\ref{sec:thue-morse}, and let
\[
 \mathcal E=\{(\theta,y):y\in Y_\theta\},\qquad
 Q=Z\times\mathcal E\times\{0,1\},\qquad
 X=Q\sqcup\{\dagger\},\qquad U=\{0,1\}^2.
\]
Hausdorff continuity makes $\mathcal E$ compact.
For $x=(z,\theta,y,r)$ define
\[
 f(x,(a,b))=
 \begin{cases}
  (\sigma z,\theta,\sigma y,b),&a=z_0,\\
  \dagger,&a\ne z_0,
 \end{cases}
 \qquad f(\dagger,u)=\dagger.
\]
These maps are continuous on clopen matching regions. At each
state exactly the two controls with $a=z_0$ are safe. Their
successors differ in the register coordinate $b$, so the four
control maps are distinct and $Q$ is controlled invariant.
The register is a state coordinate of this system; the feedback
has no additional memory or time variable.

For every legal Borel partition, the first control label of
each atom must equal $z_0$. The first state coordinate always
shifts, independently of every Borel mixture and atom refinement.
Label projection therefore gives
$\Patt_\C(F)\ge\Patt_Z(F)$ for each finite $F$.
The two-atom clopen feedback $(z_0,0)$ attains equality for
every $F$. It follows that
\[
 \inf_{\C\in\mathfrak C_{\Bor}(Q,1)}\hseq^A(\C;Q)
 =h_{\rm seq,coord}^A(Z)\quad\text{for every }A,
 \qquad M_1=\log2.
\]
The schedule in Proposition~\ref{prop:thue-morse} gives
$L_1=\log2$. A control word of length $N$ is safe exactly
when its first components match the initial $Z$ word.
Thus $r_{\rm inv}(N,Q)=|\mathcal L_N(Z)|\le8N$ and $\hinv(Q)=0$.

For $\eta\in\Theta$, use the Borel selector
\[
 s_\eta(z,\theta,y,r)=(z_0,\one_{\{\theta=\eta\}}y_0).
\]
Its four fibers are nonempty, Borel, and assigned different
controls. Its actual name set is $Z\times Y_\eta$:
the slice $\theta=\eta$ realizes the full product, and the
other slices add only the zero second sequence. Hence
\[
 \Patt_{s_\eta}(F)=\Patt_Z(F)\Patt_{Y_\eta}(F).
\]
The binary-addition argument in Proposition~\ref{prop:thue-morse}
also gives all patterns on any finite subset of
$\{4^j:j\ge0\}$: extend the prescribed digits to the largest
exponent and apply the same parity correction. Thus both factors
realize $2^n$ patterns on any $n$ positions in $B_\eta$.
It follows that $p_{s_\eta}^*(n)=4^n$ and $h^*(s_\eta)=\log4$.
Their ordinary word entropy is zero by the word bounds for
$Z$ and $Y_\eta$, and their name languages vary Hausdorff
continuously with $\eta$.

Fix $A$. By Lemma~\ref{lem:sparse-bundle}, the normalized
logarithm of the second pattern factor tends to zero almost
surely. Therefore
\[
 h_A(s_\eta)=h_{\rm seq,coord}^A(Z)\le\log2<\log4
 \quad\text{for }\nu\text{-almost every }\eta.
\]
This disproves pointwise common attainment. In the designated
family one also has
$\max_A\inf_\eta h_A(s_\eta)=\log2<\min_\eta h^*(s_\eta)=\log4$.
For the full Borel class the lower-rate selector $(z_0,0)$
is available, and the exact equality $L_1=M_1$ remains valid.
\end{proof}

The compact family here consists of name languages, in the Hausdorff
topology; no topology or compactness is imposed on the Borel selectors.
The example separates attaining every feedback's own rate from
attaining the minimum over all feedbacks.


\section*{Use of AI tools}
OpenAI ChatGPT/Codex assisted with mathematical exploration and
candidate proofs, comparison of primary literature, and manuscript
preparation, including language editing and checks of the LaTeX
source and compiled PDF.

\begin{thebibliography}{99}

\bibitem{Debauche2022}
Debauche, V., Della Rossa, M., \& Jungers, R.~M. (2022).
Comparison of path-complete Lyapunov functions via template-dependent lifts.
\emph{Nonlinear Analysis: Hybrid Systems, 46}, 101237.
\doi{10.1016/j.nahs.2022.101237}.
Author version: \href{https://arxiv.org/abs/2110.13474v1}{arXiv:2110.13474v1}.

\bibitem{Goodman1974}
Goodman, T.~N.~T. (1974). Topological sequence entropy.
\emph{Proceedings of the London Mathematical Society, 29}(2), 331--350.
\doi{10.1112/plms/s3-29.2.331}.

\bibitem{HuangMaassYe2004}
Huang, W., Maass, A., \& Ye, X. (2004).
Sequence entropy pairs and complexity pairs for a measure.
\emph{Annales de l'Institut Fourier, 54}(4), 1005--1028.
\doi{10.5802/aif.2041}.

\bibitem{HuangYe2009}
Huang, W., \& Ye, X. (2009).
Combinatorial lemmas and applications to dynamics.
\emph{Advances in Mathematics, 220}(6), 1689--1716.
\doi{10.1016/j.aim.2008.11.009}.

\bibitem{Hulse1982}
Hulse, P. (1982). Sequence entropy and subsequence generators.
\emph{Journal of the London Mathematical Society, 26}(3), 441--450.
\doi{10.1112/jlms/s2-26.3.441}.

\bibitem{KamaeZamboniSystems2002}
Kamae, T., \& Zamboni, L. (2002).
Maximal pattern complexity for discrete systems.
\emph{Ergodic Theory and Dynamical Systems, 22}(4), 1201--1214.
\doi{10.1017/S0143385702000585}.

\bibitem{Kawan2013}
Kawan, C. (2013).
\emph{Invariance entropy for deterministic control systems:
An introduction}. Lecture Notes in Mathematics, Vol.~2089.
Springer. \doi{10.1007/978-3-319-01288-9}.

\bibitem{KL2007}
Kerr, D., \& Li, H. (2007).
Independence in topological and C$^*$-dynamics.
\emph{Mathematische Annalen, 338}, 869--926.
\doi{10.1007/s00208-007-0097-z}.

\bibitem{KL2009}
Kerr, D., \& Li, H. (2009).
Combinatorial independence in measurable dynamics.
\emph{Journal of Functional Analysis, 256}(5), 1341--1386.
\doi{10.1016/j.jfa.2008.12.014}.

\bibitem{Romagnoli2007}
Romagnoli, P.-P. (2007).
A remark on the topological entropies of covers and partitions.
\emph{Studia Mathematica, 182}(3), 273--281.
\doi{10.4064/sm182-3-6}.

\bibitem{TomarKawanZamani2022}
Tomar, M.~S., Kawan, C., \& Zamani, M. (2022).
Numerical over-approximation of invariance entropy via finite abstractions.
\emph{Systems \& Control Letters, 170}, 105395.
\doi{10.1016/j.sysconle.2022.105395}.

\end{thebibliography}

\end{document}
